% ROZDZIAL 3 -- fragment do dolaczenia do glownego skryptu.
% Przykladowe programy zakladaja C++20, jesli nie podano inaczej.
% W preambule dokumentu potrzebny jest pakiet listings:
% \usepackage{listings}
% \lstset{language=C++,basicstyle=\ttfamily\small,
%   columns=fullflexible,keepspaces=true,breaklines=true,
%   showstringspaces=false,frame=single,tabsize=4}
% Fragment mozna wlaczyc przez:
% % ROZDZIAL 3 -- fragment do dolaczenia do glownego skryptu.
% Przykladowe programy zakladaja C++20, jesli nie podano inaczej.
% W preambule dokumentu potrzebny jest pakiet listings:
% \usepackage{listings}
% \lstset{language=C++,basicstyle=\ttfamily\small,
%   columns=fullflexible,keepspaces=true,breaklines=true,
%   showstringspaces=false,frame=single,tabsize=4}
% Fragment mozna wlaczyc przez:
% % ROZDZIAL 3 -- fragment do dolaczenia do glownego skryptu.
% Przykladowe programy zakladaja C++20, jesli nie podano inaczej.
% W preambule dokumentu potrzebny jest pakiet listings:
% \usepackage{listings}
% \lstset{language=C++,basicstyle=\ttfamily\small,
%   columns=fullflexible,keepspaces=true,breaklines=true,
%   showstringspaces=false,frame=single,tabsize=4}
% Fragment mozna wlaczyc przez:
% % ROZDZIAL 3 -- fragment do dolaczenia do glownego skryptu.
% Przykladowe programy zakladaja C++20, jesli nie podano inaczej.
% W preambule dokumentu potrzebny jest pakiet listings:
% \usepackage{listings}
% \lstset{language=C++,basicstyle=\ttfamily\small,
%   columns=fullflexible,keepspaces=true,breaklines=true,
%   showstringspaces=false,frame=single,tabsize=4}
% Fragment mozna wlaczyc przez:
% \input{rozdzial_3_wskazniki_smart_pointery_cpp.tex}

\chapter[Wskaźniki i inteligentne wskaźniki w C++]{Wskaźniki, zarządzanie pamięcią i inteligentne wskaźniki w C++}
\label{chap:wskazniki-cpp}

\section{Po co istnieją wskaźniki?}

Wskaźnik (ang. \emph{pointer}) przechowuje wartość pozwalającą
odwołać się do obiektu albo funkcji. Typowy wskaźnik do obiektu ma
postać \texttt{T*}. W odróżnieniu od obiektu typu \texttt{T},
wskaźnik może wskazywać na inny obiekt, może być pusty, a samo
wskazywanie \emph{nie oznacza posiadania wskazywanego obiektu}.

Wskaźniki wykorzystuje się między innymi do:
\begin{itemize}
    \item przekazywania i modyfikowania istniejących obiektów bez ich kopiowania;
    \item opisywania opcjonalnej relacji ``obiekt istnieje albo go nie ma'';
    \item implementowania dynamicznych struktur danych, takich jak listy,
    drzewa i grafy;
    \item realizowania polimorfizmu i dynamicznego wiązania;
    \item współpracy z interfejsami C, sterownikami i bibliotekami systemowymi;
    \item kontrolowania czasu życia zasobów za pomocą
    \texttt{std::unique\_ptr}, \texttt{std::shared\_ptr} i
    \texttt{std::weak\_ptr}.
\end{itemize}

\paragraph{Najważniejsza zasada.}
W nowoczesnym C++ wskaźnik \texttt{T*} traktujemy przede wszystkim
jako \emph{nieposiadający dostęp} do obiektu (\emph{observer}).
Odpowiedzialność za czas życia obiektu powinna być reprezentowana
osobno i jednoznacznie. Najczęściej wystarcza obiekt przechowywany
bezpośrednio, kontener standardowy albo \texttt{std::unique\_ptr}.

\section{Podstawy: adres, dereferencja i \texttt{nullptr}}

\begin{center}\small
\begin{tabularx}{\textwidth}{@{}p{0.32\textwidth}X@{}}
\toprule
Zapis & Za co odpowiada \\
\midrule
\texttt{\&x} & Pobiera adres obiektu \code{x}. \\
\code{*p} & Uzyskuje obiekt wskazywany przez \code{p}; wymaga ważnego, niepustego wskaźnika. \\
\code{p->field} & Udostępnia składową przez wskaźnik; odpowiada \code{(*p).field}. \\
\code{nullptr} & Oznacza brak wskazywanego obiektu. \\
\code{std::span<T>} & Łączy dostęp do kolejnych elementów z ich liczbą, ale nie posiada pamięci. \\
\code{std::addressof(x)} & Pobiera rzeczywisty adres obiektu nawet przy przeciążonym operatorze \texttt{\&}. \\
\bottomrule
\end{tabularx}
\end{center}

\subsection{Deklarowanie wskaźników}

Operator \texttt{\&} zastosowany do obiektu daje jego adres, natomiast
operator jednoargumentowy \texttt{*} służy do dereferencji (dostępu do
wskazywanego obiektu):

\begin{lstlisting}
#include <iostream>

int main() {
    int x = 10;
    int* p = &x;             // p points to x

    std::cout << x << '\n';  // 10
    std::cout << *p << '\n'; // 10

    *p = 25;                 // modifies x
    std::cout << x << '\n';  // 25
}
\end{lstlisting}

W deklaracji \texttt{int* p} gwiazdka należy do deklaratora, nie jest
operacją dereferencji. Nie należy też utożsamiać wskaźnika z samym
obiektem: \texttt{p} jest osobnym obiektem o własnej wartości i adresie.

\begin{lstlisting}
int x = 7;
int* p = &x;
int** pp = &p;  // pointer to a pointer

**pp = 9;       // changes x
\end{lstlisting}

Zapis \texttt{int* a, b;} deklaruje wskaźnik \texttt{a} oraz
zwykłą zmienną \texttt{b}. Dla czytelności lepiej deklarować je
oddzielnie.

\subsection{Wskaźnik pusty i poprawna dereferencja}

Pusty wskaźnik zapisujemy za pomocą \texttt{nullptr} (od C++11).
Nie wolno go dereferencjonować. Samo porównanie z
\texttt{nullptr} jest bezpieczne:

\begin{lstlisting}
int* p = nullptr;

if (p != nullptr) {
    *p = 5;
}

if (p) {                 // equivalent non-null test
    *p = 6;
}
\end{lstlisting}

\texttt{nullptr} ma osobny typ \texttt{std::nullptr\_t}; w odróżnieniu od
historycznych zapisów \texttt{0} i \texttt{NULL} nie jest zwykłą
liczbą całkowitą. Pomaga to przy przeciążaniu funkcji:

\begin{lstlisting}
void f(int);
void f(int*);

f(0);         // calls f(int)
f(nullptr);   // calls f(int*)
\end{lstlisting}

\paragraph{Uwaga o wskaźnikach niezainicjalizowanych.}
Lokalne \texttt{int* p;} bez inicjalizacji ma nieokreśloną wartość;
korzystanie z niej może prowadzić do niezdefiniowanego zachowania.
Zawsze inicjalizuj wskaźnik poprawnym adresem lub \texttt{nullptr}.

\subsection{Dostęp do składowych: \texttt{.} i \texttt{->}}

Operator \texttt{.} wybiera składową obiektu, natomiast
\texttt{->} wybiera składową przez wskaźnik:

\begin{lstlisting}
#include <string>

struct User {
    std::string name;
    int score = 0;
};

User u{"Anna", 12};
User* p = &u;

u.score = 13;
p->score = 14;       // same as (*p).score = 14
\end{lstlisting}

Nawiasy w \texttt{(*p).score} są potrzebne ze względu na priorytety
operatorów. Zapis \texttt{*p.score} nie oznacza tego samego.

\section{Wskaźniki, referencje i kwalifikator \texttt{const}}

\subsection{Kiedy wskaźnik, a kiedy referencja?}

Referencja \texttt{T\&} jest zwykle odpowiednia, gdy obiekt musi
istnieć i parametr nie jest opcjonalny. Wskaźnik \texttt{T*} nadaje
się, gdy brak obiektu jest dopuszczalny lub gdy chcemy zmieniać
obiekt, na który wskazujemy.

\begin{lstlisting}
void increase(int& x) { ++x; }

void increase_if_present(int* x) {
    if (x) { ++*x; }
}

int n = 3;
increase(n);             // n == 4
increase_if_present(&n); // n == 5
increase_if_present(nullptr); // no operation
\end{lstlisting}

Referencja również może stać się \emph{wisząca}, gdy obiekt,
do którego się odnosi, przestanie istnieć. Typ \texttt{T\&}
nie gwarantuje samodzielnie, że czas życia obiektu jest wystarczająco długi.

\subsection{Co jest stałe: wskaźnik czy obiekt?}

\begin{lstlisting}
int a = 1;
int b = 2;

const int* p = &a;       // read-only access through p
p = &b;                 // OK: p can point elsewhere
// *p = 3;              // error

int* const q = &a;       // q cannot be reassigned
*q = 3;                 // OK: a is mutable
// q = &b;              // error

const int* const r = &a; // neither reassignment nor write
int const* s = &a;       // same type as const int*
\end{lstlisting}

Reguła praktyczna: \texttt{const} po \texttt{*} dotyczy samego
wskaźnika, a przed \texttt{*} --- wskazywanego typu.

\paragraph{Pozorna niemutowalność.}
Wskaźnik \texttt{const T*} nie pozwala modyfikować obiektu
\emph{przez ten wskaźnik}. Nie oznacza, że obiekt nie może ulec zmianie
inną drogą:

\begin{lstlisting}
int x = 10;
const int* read_only = &x;
x = 20;                 // allowed
// *read_only = 30;     // not allowed
\end{lstlisting}

\section{Czas życia obiektu i wskaźniki wiszące}

\subsection{Adres nie oznacza prawa dostępu}

Wskaźnik może przechowywać adres miejsca, w którym obiekt już
nie istnieje. Jest wtedy wskaźnikiem \emph{wiszącym}
(ang. \emph{dangling pointer}). Dereferencja takiego wskaźnika
prowadzi do niezdefiniowanego zachowania (\emph{undefined behavior}, UB).

\begin{lstlisting}
int* bad() {
    int local = 42;
    return &local;       // WRONG: local dies at return
}

// int x = *bad();       // UB: dangling pointer
\end{lstlisting}

Prawidłowo można zwrócić samą wartość:

\begin{lstlisting}
int good() {
    int local = 42;
    return local;        // safe value return
}
\end{lstlisting}

\subsection{Unieważnianie wskaźników przez kontenery}

Nie każdy wskaźnik staje się wiszący przez wyjście z funkcji.
Relokacja elementów \texttt{std::vector} także może unieważnić
wskaźniki, referencje i iteratory do tych elementów:

\begin{lstlisting}
#include <vector>

std::vector<int> values{10, 20};
int* p = &values[0];

values.reserve(100);     // may reallocate storage
// *p = 50;              // unsafe if reallocation occurred
p = &values[0];          // reacquire a valid pointer
\end{lstlisting}

Po zmianie pojemności wymagającej realokacji adresy elementów
się zmieniają. W ogólności po operacji mogącej unieważnić wskaźnik
należy pobrać go ponownie, a nie próbować odgadnąć, czy nadal działa.

\subsection{Trzy niezależne pytania}

Przy każdym wskaźniku rozstrzygaj:
\begin{enumerate}
    \item Czy wartość wskaźnika jest poprawna (np. nie jest
    niezainicjalizowana)?
    \item Czy obiekt, na który wskazujemy, nadal istnieje?
    \item Czy operacja jest dozwolona dla danego typu i zakresu pamięci
    (np. czy nie wychodzi poza tablicę)?
\end{enumerate}

Sam test \texttt{p != nullptr} rozwiązuje co najwyżej fragment
pierwszego problemu, ale \emph{nie} wykrywa wiszących wskaźników.

\section{Tablice, arytmetyka wskaźników i \texttt{std::span}}

\subsection{Arytmetyka wskaźników}

Dla wskaźnika do elementu tablicy operacje \texttt{p + k} i
\texttt{p[k]} odnoszą się do jej elementów, pod warunkiem że
poruszamy się w dozwolonym zakresie. Można utworzyć wskaźnik
\emph{jeden za końcem} tablicy, ale nie wolno go dereferencjonować.

\begin{lstlisting}
int a[]{4, 8, 12};
int* p = a;          // points to a[0]

int x = *(p + 1);    // 8
int y = p[2];       // 12
int* end = a + 3;   // valid one-past-the-end pointer
// int z = *end;     // UB
\end{lstlisting}

Dodawanie i odejmowanie wskaźników ma sens w obrębie \emph{tej samej}
tablicy (z uwzględnieniem pozycji jeden za końcem). Nie należy
wykorzystywać takich działań do przemieszczania się między dowolnymi
obiektami w pamięci.

\subsection{Tablica to nie wskaźnik}

Typ \texttt{int[3]} i typ \texttt{int*} są różne, choć w wielu
wyrażeniach tablica jest automatycznie konwertowana na wskaźnik
do pierwszego elementu (tzw. \emph{array-to-pointer decay}).

\begin{lstlisting}
int a[3] = {1, 2, 3};
int* p = a;

static_assert(sizeof(a) == 3 * sizeof(int));
// sizeof(p) is the size of a pointer, not of the array
\end{lstlisting}

Funkcja \texttt{void f(int* p)} nie zna długości tablicy.
Jeśli trzeba pracować z ciągłym fragmentem pamięci o znanym
rozmiarze, C++20 udostępnia \texttt{std::span}:

\begin{lstlisting}
#include <span>
#include <vector>

int sum(std::span<const int> xs) {
    int result = 0;
    for (int x : xs) {
        result += x;
    }
    return result;
}

int a[]{1, 2, 3};
std::vector<int> v{4, 5, 6};

int s1 = sum(a);  // 6
int s2 = sum(v);  // 15
\end{lstlisting}

\texttt{std::span<T>} przechowuje widok ciągłych danych
(z informacją o liczbie elementów), ale \emph{nie jest właścicielem}
pamięci. Gdy dane znikną albo przeniosą się w inne miejsce,
\texttt{span} również może stać się niepoprawny. Dla większości
dynamicznych tablic własnościowych preferuj \texttt{std::vector<T>}.

\section{\texttt{void*}, konwersje i rzutowania wskaźników}

\subsection{\texttt{void*}: adres bez informacji o typie}

Wskaźnik \texttt{void*} może przechowywać adres obiektu bez
informacji o jego konkretnym typie. Bywa używany w interfejsach C
i bibliotekach niskopoziomowych, ale nie pozwala na bezpośrednią
dereferencję:

\begin{lstlisting}
int x = 42;
void* opaque = &x;              // type erased
// *opaque = 7;                 // error: cannot dereference void*

int* p = static_cast<int*>(opaque);
*p = 7;                        // valid: original object is int
\end{lstlisting}

Przywrócenie błędnego typu i późniejsza próba użycia obiektu
mogą prowadzić do niezdefiniowanego zachowania. Jeśli typ jest
znany, zwykle lepiej przechowywać od razu \texttt{T*}.

\subsection{\texttt{static\_cast}, \texttt{dynamic\_cast} i
\texttt{reinterpret\_cast}}

\begin{itemize}
    \item \texttt{static\_cast<T*>} nadaje się m.in. do poprawnych
    konwersji pomiędzy \texttt{void*} a wskaźnikiem obiektowym,
    jeśli rzeczywisty typ i wyrównanie na to pozwalają.
    \item \texttt{dynamic\_cast<T*>} służy do sprawdzanych w czasie
    wykonania konwersji w hierarchiach klas polimorficznych.
    Niepowodzenie konwersji wskaźnika zwraca \texttt{nullptr}.
    \item \texttt{reinterpret\_cast<T*>} umożliwia niskopoziomową
    reinterpretację reprezentacji wskaźnika, ale sama konwersja
    nie uprawnia do dowolnej dereferencji ani nie tworzy obiektu
    odpowiedniego typu.
\end{itemize}

\begin{lstlisting}
struct Base {
    virtual ~Base() = default;
};

struct Derived : Base {
    int value = 12;
};

Derived d;
Base* base = &d;

if (Derived* p = dynamic_cast<Derived*>(base)) {
    p->value = 30;          // checked downcast succeeded
}
\end{lstlisting}

Przy przekazywaniu parametrów funkcji nie używaj rzutowań
jako sposobu na obejście \texttt{const}, sprawdzania czasu życia
ani zasad własności. Dobrze zaprojektowany interfejs zazwyczaj
w ogóle nie wymaga rzutowania wskaźników.

\section{Pamięć dynamiczna: \texttt{new}, \texttt{delete} i ich ograniczenia}

\subsection{Klasyczny mechanizm}

Operator \texttt{new} tworzy obiekt w pamięci przydzielonej
dynamicznie i zwraca wskaźnik. Odpowiadający mu
\texttt{delete} niszczy obiekt i zwalnia pamięć:

\begin{lstlisting}
int* p = new int(42);
// use *p

delete p;
p = nullptr;           // avoids this one stale alias
\end{lstlisting}

Dla tablic obowiązuje para \texttt{new[]} i \texttt{delete[]}:

\begin{lstlisting}
int* a = new int[5]{}; // zero-initialized elements
// use a[0], ..., a[4]

delete[] a;
\end{lstlisting}

\textbf{Nie wolno} łączyć \texttt{new[]} z \texttt{delete} ani
zwykłego \texttt{new} z \texttt{delete[]}. Ponowne
\texttt{delete} tego samego żywego przydziału również jest błędem.
\texttt{delete nullptr;} jest natomiast bezpieczne.

\subsection{Dlaczego ręczne zarządzanie pamięcią jest niebezpieczne?}

\begin{lstlisting}
void fragile() {
    int* p = new int(42);
    operation_that_may_throw();
    delete p;   // skipped if the operation throws
}
\end{lstlisting}

Jeśli funkcja zgłosi wyjątek, pamięć zostaje utracona
(\emph{memory leak}). Analogiczny problem powstaje przy
wielu instrukcjach \texttt{return}, kopiowaniu wskaźników i
niespójnym przekazywaniu odpowiedzialności za \texttt{delete}.

\paragraph{Zalecenie.}
W typowym kodzie aplikacyjnym nie pisz bezpośrednio
\texttt{new}/\texttt{delete}. Użyj obiektu lokalnego,
\texttt{std::vector<T>} albo inteligentnego wskaźnika.
Ręczne operacje alokowania pozostają potrzebne w implementacjach
niskopoziomowych, alokatorach i przy integracji z określonym API.

\section{RAII i semantyka własności}

\subsection{RAII: zasób należy do obiektu}

\emph{Resource Acquisition Is Initialization} (RAII) oznacza
powiązanie przejęcia zasobu z utworzeniem obiektu zarządzającego,
a zwolnienia zasobu z jego destruktorem. Stosuje się je nie tylko
do pamięci, ale też do plików, blokad, gniazd sieciowych i uchwytów
systemowych.

\begin{lstlisting}
#include <fstream>
#include <string>

void save(const std::string& path) {
    std::ofstream file(path);
    file << "example\n";
} // file is closed by its destructor
\end{lstlisting}

Destruktor uruchamia się podczas normalnego opuszczania zakresu
oraz podczas odwijania stosu po wyjątku. Stąd RAII upraszcza
bezpieczeństwo wyjątkowe.

\subsection{Trzy rodzaje relacji z obiektem}

\begin{description}
    \item[Własność wyłączna (\emph{unique ownership}).]
    Dokładnie jeden obiekt zarządzający odpowiada za zniszczenie
    zasobu: najczęściej \texttt{std::unique\_ptr<T>}.
    \item[Własność współdzielona (\emph{shared ownership}).]
    Wielu właścicieli może utrzymywać jeden zasób przy życiu:
    \texttt{std::shared\_ptr<T>}.
    \item[Dostęp bez własności (\emph{non-owning access}).]
    Użytkownik widzi obiekt, ale nie przedłuża jego życia:
    zwykły wskaźnik, referencja, \texttt{std::span<T>} lub
    \texttt{std::weak\_ptr<T>} (ten ostatni obserwuje stan
    zarządzany przez \texttt{shared\_ptr} i pozwala warunkowo
    uzyskać właściciela).
\end{description}

Dobór reprezentacji wynika z \emph{kontraktu własności}, a nie
z samego faktu, że obiekt znajduje się w pamięci dynamicznej.

\section{\texttt{std::unique\_ptr}: jeden właściciel}

\begin{center}\small
\begin{tabularx}{\textwidth}{@{}p{0.36\textwidth}X@{}}
\toprule
Mechanizm & Własność i czas życia \\
\midrule
\code{std::make_unique<T>(...)} & Tworzy obiekt i zwraca jego jedynego właściciela. \\
\code{unique_ptr::get()} & Zwraca surowy wskaźnik bez przekazania własności. \\
\code{unique_ptr::reset(p)} & Niszczy obecnie posiadany obiekt i przejmuje wskaźnik \code{p}. \\
\code{unique_ptr::release()} & Oddaje surowy wskaźnik bez zniszczenia obiektu; odtąd trzeba osobno zadbać o jego zwolnienie. \\
\code{std::shared_ptr<T>} & Współdzieli własność; obiekt znika po odejściu ostatniego właściciela. \\
\code{weak_ptr::lock()} & Próbuje uzyskać \code{shared_ptr}; zwraca pusty wynik, jeśli obiekt wygasł. \\
\bottomrule
\end{tabularx}
\end{center}

\subsection{Tworzenie obiektu i automatyczne niszczenie}

\texttt{std::unique\_ptr<T>} (nagłówek \texttt{<memory>}, C++11)
zarządza pojedynczym obiektem przez wyłączną własność. Od C++14
najczęściej tworzymy go przez \texttt{std::make\_unique<T>}:

\begin{lstlisting}
#include <memory>
#include <string>

#include <utility>

struct Person {
    explicit Person(std::string n) : name(std::move(n)) {}
    std::string name;
};

void example() {
    auto p = std::make_unique<Person>("Ala");
    p->name = "Ola";
    // no delete needed
} // destructor destroys Person
\end{lstlisting}

\texttt{unique\_ptr} można traktować jako obiekt zarządzający:
jego destruktor wywołuje odpowiednią funkcję usuwającą zasób
(\emph{deleter}). Domyślnie jest nią \texttt{std::default\_delete<T>}.

\subsection{Brak kopiowania i przenoszenie własności}

Wskaźnik o własności wyłącznej nie może być kopiowany.
Może natomiast zostać \emph{przeniesiony} za pomocą
\texttt{std::move}:

\begin{lstlisting}
#include <memory>
#include <utility>

auto first = std::make_unique<int>(10);
// auto bad = first;                  // error: copy

auto second = std::move(first);      // transfers ownership
// first == nullptr
*second = 20;
\end{lstlisting}

\texttt{std::move} samo niczego nie zwalnia i nie przenosi fizycznie
wskazywanego obiektu. Zwraca wyrażenie pozwalające wywołać konstruktor
lub przypisanie przenoszące. W tym przypadku przenoszona jest
\emph{własność wskaźnika}, a obiekt typu \texttt{int} nie zmienia adresu.

\subsection{Przekazywanie do funkcji: kto jest właścicielem?}

Rozróżnijmy trzy interfejsy:

\begin{lstlisting}
#include <memory>
#include <utility>

struct Widget { void run() {} };

void inspect(const Widget& w);             // borrows, non-null
void maybe_inspect(const Widget* w);       // borrows, nullable
void take(std::unique_ptr<Widget> w);      // takes ownership

void demo() {
    auto p = std::make_unique<Widget>();
    inspect(*p);                // p still owns the object
    maybe_inspect(p.get());     // p still owns the object
    take(std::move(p));         // ownership transferred
    // p is now empty
}
\end{lstlisting}

Jeżeli funkcja \emph{jedynie używa} obiektu, przekazywanie przez
\texttt{T\&}, \texttt{const T\&} lub \texttt{T*} jest zwykle
czytelniejsze niż przekazywanie \texttt{unique\_ptr} przez referencję.
Parametr \texttt{unique\_ptr<T>} przyjmowany \emph{przez wartość}
sygnalizuje przejęcie własności. Dla funkcji, która ma zastąpić
obiekt będący własnością wywołującego, można użyć
\texttt{std::unique\_ptr<T>\&}.

\subsection{Metody \texttt{get}, \texttt{reset}, \texttt{release}, \texttt{swap}}

\begin{lstlisting}
#include <memory>
#include <utility>

auto p = std::make_unique<int>(7);
int* observer = p.get();       // borrow; do not delete

p.reset(new int(9));           // destroys old int
// observer is now dangling

auto q = std::make_unique<int>(11);
p.swap(q);                    // exchanges managed objects

int* raw = q.release();        // q becomes empty; raw owns now
// responsibility for delete raw is now explicit
delete raw;
\end{lstlisting}

\texttt{get()} zwraca dostęp bez własności. \texttt{reset()} zastępuje
zarządzany obiekt, zwalniając dotychczasowy. \texttt{release()} oddaje
surowy wskaźnik \emph{bez zwolnienia obiektu}; dlatego jest potrzebne
głównie przy przekazywaniu własności do interfejsu, który tego wymaga.
Nie należy mylić \texttt{release()} z \texttt{reset()}.

\subsection{Tablice: \texttt{unique\_ptr<T[]>} a \texttt{vector<T>}}

\begin{lstlisting}
#include <memory>
#include <vector>

auto a = std::make_unique<int[]>(5); // array of 5 zeros
a[2] = 17;

std::vector<int> b(5);         // usually more convenient
b.at(2) = 17;                  // checked indexing
\end{lstlisting}

\texttt{unique\_ptr<T[]>} dobiera \texttt{delete[]} i udostępnia
\texttt{operator[]}, lecz sam nie przechowuje rozmiaru.
Zwykle preferuj \texttt{std::vector<T>}, bo posiada rozmiar,
iteratory i bogatszy interfejs. Dla tablicy o stałym rozmiarze
określonym w kodzie rozważ \texttt{std::array<T, N>}.

\subsection{Własny deleter: zasób inny niż pamięć \texttt{new}}

Inteligentny wskaźnik może wywoływać dowolną funkcję zwalniającą
zasób. Typowym przykładem jest plik w interfejsie C:

\begin{lstlisting}
#include <cstdio>
#include <memory>

struct FileCloser {
    void operator()(std::FILE* file) const noexcept {
        if (file) { std::fclose(file); }
    }
};

using FilePtr = std::unique_ptr<std::FILE, FileCloser>;

void example() {
    FilePtr file(std::fopen("log.txt", "w"));
    if (!file) { return; }

    std::fputs("hello\n", file.get());
} // fclose is called automatically
\end{lstlisting}

W praktyce deleter bywa osobną strukturą, a jego typ wchodzi
w skład typu \texttt{unique\_ptr<T, Deleter>}. Mechanizm pozwala
owinąć RAII również zasoby zwalniane przez \texttt{free()},
\texttt{fclose()} lub bibliotekę zewnętrzną.

\subsection{Polimorfizm i wirtualny destruktor}

\begin{lstlisting}
#include <memory>

struct Shape {
    virtual ~Shape() = default;
    virtual double area() const = 0;
};

struct Circle : Shape {
    explicit Circle(double r) : radius(r) {}
    double area() const override { return 3.14159 * radius * radius; }
    double radius;
};

std::unique_ptr<Shape> make_shape() {
    return std::make_unique<Circle>(2.0);
}
\end{lstlisting}

Przy niszczeniu obiektu pochodnego przez wskaźnik do bazy
zwykle potrzebny jest wirtualny destruktor bazy.
Dotyczy to również przenoszenia własności z klasy pochodnej
do bazowej przy użyciu \texttt{unique\_ptr} i domyślnego deletera.

\section{\texttt{std::shared\_ptr}: współdzielona własność}

\subsection{Jak działa licznik właścicieli?}

\texttt{std::shared\_ptr<T>} (C++11) udostępnia współdzieloną
własność zasobu. Zwykle tworzymy go za pomocą
\texttt{std::make\_shared<T>}:

\begin{lstlisting}
#include <memory>

auto a = std::make_shared<int>(10);
auto b = a;             // shared ownership

// a.use_count() == 2 (provided there are no other owners)
b.reset();              // releases b's ownership
// a.use_count() == 1
\end{lstlisting}

Gdy liczba \emph{silnych właścicieli} spadnie do zera, zarządzany
obiekt zostaje zniszczony. Implementacja utrzymuje tzw.
\emph{blok kontrolny} (\emph{control block}) zawierający między
innymi liczniki właścicieli, informacje potrzebne do zniszczenia
zasobu i ewentualny deleter.

\paragraph{Schemat ideowy.}
\begin{lstlisting}
shared_ptr a ---\
                +--> [ control block: strong = 2 ] --> [ T ]
shared_ptr b ---/           weak = 0
\end{lstlisting}

\texttt{use\_count()} ma znaczenie przede wszystkim diagnostyczne.
Nie powinno się na podstawie odczytanego licznika podejmować decyzji
o bezpieczeństwie wątkowym ani o wyłącznej własności, która może
zmienić się równolegle.

\subsection{Tworzenie i przekazywanie \texttt{shared\_ptr}}

\begin{lstlisting}
#include <memory>

struct Data { int x = 0; };

std::shared_ptr<Data> create() {
    return std::make_shared<Data>();
}

void keep(std::shared_ptr<Data> p) {
    // stores a copy of p somewhere, if needed
}

void read_only(const Data& data) {
    // does not require shared ownership
}

void demo() {
    auto p = create();
    read_only(*p); // borrow
    keep(p);       // creates another temporary shared owner
}
\end{lstlisting}

Parametr \texttt{shared\_ptr<T>} przyjmowany przez wartość oznacza,
że funkcja otrzymuje własny udział we własności na czas wywołania
i ewentualnie może go przechować. Jeśli funkcja wyłącznie
czyta obiekt, zwykle wystarcza \texttt{const T\&} lub \texttt{T*}.
Samo przekazanie \texttt{const shared\_ptr<T>\&} unika inkrementacji
licznika, ale bardziej wiąże interfejs z modelem własności.

\subsection{\texttt{make\_shared} a konstruktor ze wskaźnika}

W typowych implementacjach \texttt{make\_shared} umieszcza obiekt
i blok kontrolny w jednym przydziale pamięci, co może ograniczać
narzut alokacji. Natomiast konstrukcja \texttt{shared\_ptr<T>(new T)}
zwykle wymaga osobnych przydziałów.

\begin{lstlisting}
auto preferred = std::make_shared<int>(42);
std::shared_ptr<int> also_valid(new int(42));
\end{lstlisting}

Nie wolno jednak tworzyć dwóch \texttt{shared\_ptr} niezależnie
z \emph{tego samego} surowego wskaźnika:

\begin{lstlisting}
int* raw = new int(42);
std::shared_ptr<int> a(raw);
// std::shared_ptr<int> b(raw); // WRONG: second control block
\end{lstlisting}

Dwa niezależne bloki kontrolne próbowałyby zwolnić ten sam obiekt.
Poprawne współdzielenie powstaje przez \emph{kopiowanie istniejącego}
\texttt{shared\_ptr}: \texttt{auto b = a;}.

\paragraph{Kiedy nie stosować \texttt{make\_shared}?}
Zintegrowany przydział może zatrzymywać pamięć zajmowaną przez
obiekt w przydzielonym bloku do momentu zniszczenia także ostatniego
\texttt{weak\_ptr}, nawet gdy destruktor obiektu już się wykonał.
Przy dużych obiektach i długo żyjących słabych referencjach
warto rozważyć osobny przydział. Innymi powodami mogą być
specjalne wymagania alokacji lub konstruktory niedostępne
dla \texttt{make\_shared}.

\subsection{Współdzielenie własności a współdzielenie mutacji}

Sam \texttt{shared\_ptr} nie synchronizuje dostępu do pól obiektu:

\begin{lstlisting}
#include <memory>
#include <mutex>

struct Counter {
    std::mutex mutex;
    int value = 0;

    void increment() {
        std::lock_guard<std::mutex> lock(mutex);
        ++value;
    }
};

auto p = std::make_shared<Counter>();
// many threads may own copies of p, but must synchronize
// concurrent modifications of Counter::value
\end{lstlisting}

Blok kontrolny pozwala różnym egzemplarzom \texttt{shared\_ptr}
bezpiecznie współdzielić własność. Równoczesna modyfikacja
\emph{tego samego} egzemplarza wymaga jednak synchronizacji.
Od C++20 można użyć \texttt{std::atomic<std::shared\_ptr<T>>}.

\section{\texttt{std::weak\_ptr}: obserwacja bez utrzymywania przy życiu}

\subsection{Problem cykli własności}

Gdy obiekty wzajemnie posiadają się przez \texttt{shared\_ptr},
liczniki właścicieli mogą nigdy nie spaść do zera:

\begin{lstlisting}
#include <memory>

struct Node {
    std::shared_ptr<Node> next;
};

void leak_cycle() {
    auto a = std::make_shared<Node>();
    auto b = std::make_shared<Node>();
    a->next = b;
    b->next = a;          // cycle of strong ownership
} // nodes remain alive after both locals are destroyed
\end{lstlisting}

To nie jest problem \emph{niezerowego wskaźnika}, lecz błędnego
modelu własności: każdy węzeł utrzymuje przy życiu drugi.

\subsection{Przerywanie cyklu i \texttt{lock()}}

\texttt{std::weak\_ptr<T>} obserwuje ten sam blok kontrolny,
ale nie zwiększa liczby silnych właścicieli:

\begin{lstlisting}
#include <memory>

struct Node {
    std::weak_ptr<Node> parent;         // non-owning
    std::shared_ptr<Node> child;        // owning
};

void example() {
    auto root = std::make_shared<Node>();
    auto leaf = std::make_shared<Node>();

    root->child = leaf;
    leaf->parent = root;

    if (auto parent = leaf->parent.lock()) {
        // parent is a temporary shared owner
        (void)parent;
    }
} // no ownership cycle
\end{lstlisting}

Metoda \texttt{lock()} atomowo próbuje uzyskać
\texttt{shared\_ptr<T>}. Jeżeli obiekt jeszcze istnieje,
zwraca właściciela utrzymującego go przy życiu przez czas
używania. W przeciwnym razie zwraca pusty \texttt{shared\_ptr}.

\begin{lstlisting}
std::weak_ptr<int> observer;
{
    auto p = std::make_shared<int>(8);
    observer = p;
    if (auto alive = observer.lock()) {
        *alive += 1;           // safe while alive owns it
    }
}

// observer.expired() is now true
// observer.lock() now returns an empty shared_ptr
\end{lstlisting}

W kodzie współbieżnym nie stosuj schematu
\texttt{if (!w.expired()) { ... }} z późniejszym odwołaniem:
między sprawdzeniem a użyciem obiekt może zniknąć.
Preferuj jedną operację \texttt{lock()} i posługuj się
otrzymanym \texttt{shared\_ptr}.

\section{\texttt{enable\_shared\_from\_this} i wskaźnik \texttt{this}}

Gdy metoda obiektu zarządzanego przez \texttt{shared\_ptr}
chce przekazać dalej \emph{współdzieloną własność do siebie},
nie może wykonać \texttt{std::shared\_ptr<T>(this)}:
powstałby drugi, niezależny blok kontrolny.

Można zamiast tego dziedziczyć publicznie po
\texttt{std::enable\_shared\_from\_this<T>}:

\begin{lstlisting}
#include <memory>

struct Session : std::enable_shared_from_this<Session> {
    std::shared_ptr<Session> self() {
        return shared_from_this();
    }
};

void example() {
    auto session = std::make_shared<Session>();
    auto second_owner = session->self();
    // both share the same control block
}
\end{lstlisting}

Metodę \texttt{shared\_from\_this()} wolno wywołać dopiero wtedy,
gdy obiektem zarządza już \texttt{shared\_ptr}. Jej wywołanie
na obiekcie lokalnym albo w konstruktorze, zanim powstanie
właściciel, zgłasza wyjątek \texttt{std::bad\_weak\_ptr}.

Od C++17 dostępna jest również metoda \texttt{weak\_from\_this()}.

\section{Zaawansowany przykład: wskaźnik aliasujący \texttt{shared\_ptr}}

\texttt{shared\_ptr} rozróżnia \emph{obiekt zarządzany} (którego
czas życia kontroluje blok kontrolny) i \emph{przechowywany
wskaźnik} (zwracany przez \texttt{get()}). Specjalny konstruktor
aliasujący pozwala udostępnić adres składowej, jednocześnie
utrzymując przy życiu cały obiekt:

\begin{lstlisting}
#include <memory>

struct Pair {
    int first = 1;
    int second = 2;
};

void example() {
    auto whole = std::make_shared<Pair>();
    std::shared_ptr<int> member(whole, &whole->second);

    whole.reset();       // Pair is still alive
    *member = 20;        // accesses Pair::second safely
} // member destroys the last strong owner
\end{lstlisting}

Konstruktor aliasujący jest przydatny np. przy udostępnianiu
podobiektów, ale wymaga prawidłowej relacji czasów życia.
Nie należy go traktować jako sposobu na naprawienie
wcześniej wiszącego wskaźnika.

\section{Wskaźniki w interfejsach funkcji i klas}

\subsection{Katalog typowych kontraktów}

\begin{center}
\small
\begin{tabular}{|p{0.31\textwidth}|p{0.56\textwidth}|}
\hline
\textbf{Interfejs} & \textbf{Znaczenie} \\
\hline
\texttt{void f(const T\&)} & Obiekt wymagany, tylko odczyt; brak przekazania własności. \\
\hline
\texttt{void f(T\&)} & Obiekt wymagany, możliwa modyfikacja; brak przekazania własności. \\
\hline
\texttt{void f(const T*)} & Obiekt opcjonalny, tylko odczyt; brak własności. \\
\hline
\texttt{void f(T*)} & Obiekt opcjonalny, możliwa modyfikacja; brak własności. \\
\hline
\texttt{void f(unique\_ptr<T>)} & Przejęcie wyłącznej własności. \\
\hline
\texttt{unique\_ptr<T> f()} & Przekazanie wyłącznej własności wywołującemu. \\
\hline
\texttt{void f(shared\_ptr<T>)} & Uzyskanie udziału we współdzielonej własności. \\
\hline
\texttt{void f(weak\_ptr<T>)} & Możliwość obserwacji obiektu, którego życie kontrolują inni. \\
\hline
\texttt{void f(span<const T>)} & Widok na ciągłą sekwencję, bez własności. \\
\hline
\end{tabular}
\end{center}

W tabeli pominięto przestrzeń nazw \texttt{std::} dla czytelności.
To \emph{konwencje projektowania API}, a nie automatyczne
gwarancje długowieczności surowych wskaźników i referencji.

\subsection{Dlaczego nie zawsze używać \texttt{shared\_ptr}?}

Rozważmy funkcję, która wyświetla nazwę:

\begin{lstlisting}
#include <iostream>
#include <memory>
#include <string>

struct Product { std::string name; };

void print(const Product& p) {
    std::cout << p.name << '\n';
}

void example() {
    Product local{"book"};
    print(local);                  // no allocation

    auto owned = std::make_unique<Product>(Product{"pen"});
    print(*owned);                 // still no extra owner
}
\end{lstlisting}

Funkcja \texttt{print} nie potrzebuje znać strategii
przechowywania obiektu. Przekazanie referencji oddziela
\emph{sposób dostępu} od \emph{sposobu zarządzania czasem życia}.

\section{PImpl: ukrywanie implementacji przez \texttt{unique\_ptr}}

Technika \emph{pointer to implementation} pozwala utrzymać
szczegóły implementacji poza publicznym nagłówkiem.
Typ implementacyjny może być niekompletny w nagłówku, ale
miejsce wykonania domyślnego deletera musi widzieć jego definicję.
Dlatego destruktor klasy definiujemy w pliku \texttt{.cpp}.

\begin{lstlisting}
// widget.hpp
#pragma once
#include <memory>

class Widget {
public:
    Widget();
    ~Widget();                 // defined in .cpp

    Widget(Widget&&) noexcept;
    Widget& operator=(Widget&&) noexcept;

    Widget(const Widget&) = delete;
    Widget& operator=(const Widget&) = delete;

    void work();

private:
    struct Impl;
    std::unique_ptr<Impl> impl_;
};
\end{lstlisting}

\begin{lstlisting}
// widget.cpp
#include "widget.hpp"
#include <memory>

struct Widget::Impl {
    int value = 0;
    void run() { ++value; }
};

Widget::Widget() : impl_(std::make_unique<Impl>()) {}
Widget::~Widget() = default;
Widget::Widget(Widget&&) noexcept = default;
Widget& Widget::operator=(Widget&&) noexcept = default;

void Widget::work() { impl_->run(); }
\end{lstlisting}

Ten przykład prezentuje semantykę przenoszenia. Po przeniesieniu
źródłowy obiekt \texttt{Widget} ma pusty \texttt{impl\_};
nie wolno wywoływać na nim \texttt{work()} bez odpowiedniego
sprawdzenia. W produkcyjnym API można ten stan jawnie obsłużyć
lub określić w kontrakcie.

\section{Wskaźniki do funkcji i składowych klas}

Wskaźniki nie muszą odnosić się do danych. Można wskazywać
funkcje, a także składowe klas, choć te ostatnie wymagają
innej składni.

\subsection{Wskaźnik do funkcji}

\begin{lstlisting}
int add(int a, int b) { return a + b; }
int mul(int a, int b) { return a * b; }

using Operation = int(*)(int, int);

int apply(int x, int y, Operation op) {
    return op(x, y);
}

int a = apply(2, 3, &add); // 5
int b = apply(2, 3, &mul); // 6
\end{lstlisting}

Obiekty wywoływalne różnych typów można przechowywać za pomocą
\texttt{std::function}. Prosty szablon lub lambda bywają jednak
bardziej efektywne.

\subsection{Wskaźnik do składowej}

\begin{lstlisting}
struct Record {
    int value = 0;
    int twice() const { return 2 * value; }
};

int Record::* field = &Record::value;
int (Record::* method)() const = &Record::twice;

Record r{7};
int x = r.*field;          // 7
int y = (r.*method)();     // 14

Record* p = &r;
int z = p->*field;         // 7
\end{lstlisting}

Wskaźnik do składowej nie jest zwykłym adresem gotowego obiektu.
Dopiero w połączeniu z konkretną instancją klasy pozwala
odwołać się do odpowiedniej składowej.

\section{Wskaźniki i wielowątkowość}

\subsection{\texttt{shared\_ptr} nie eliminuje wyścigów danych}

Możliwość bezpiecznego utrzymania czasu życia obiektu nie oznacza,
że jego pola można dowolnie zmieniać z wielu wątków.
Jeśli dwa wątki zapisują do \texttt{*p}, potrzebna jest osobna
synchronizacja (np. \texttt{std::mutex} albo atomiki).

Od C++20 dostępna jest specjalizacja
\texttt{std::atomic<std::shared\_ptr<T>>}, przydatna do atomowego
publikowania nowej wersji współdzielonego obiektu:

\begin{lstlisting}
#include <atomic>
#include <memory>

struct Config {
    int workers = 4;
};

std::atomic<std::shared_ptr<const Config>> current{
    std::make_shared<const Config>(Config{4})
};

void publish(int count) {
    current.store(std::make_shared<const Config>(Config{count}));
}

int read_workers() {
    auto snapshot = current.load(); // keeps snapshot alive
    return snapshot->workers;
}
\end{lstlisting}

Tu czytelnicy otrzymują stabilny, niemutowalny \emph{snapshot}.
Atomowość dotyczy odczytu i zapisu \texttt{shared\_ptr},
a nie arbitralnych operacji na jego zawartości.

\section{Najczęstsze błędy zebrane w jednym miejscu}

\begin{enumerate}
    \item \textbf{Dereferencja \texttt{nullptr} lub wskaźnika
    niezainicjalizowanego.} Zawsze sprawdzaj kontrakt i inicjalizuj
    wskaźniki.
    \item \textbf{Wskaźnik wiszący.} \texttt{get()}, \texttt{span}
    i surowa referencja nie przedłużają życia zasobu.
    \item \textbf{Podwójne \texttt{delete}.} Jeden przydział musi mieć
    jedną, jednoznaczną strategię zwalniania.
    \item \textbf{Mieszanie \texttt{new[]} i \texttt{delete}.}
    Do tablic służy \texttt{delete[]}; zwykle lepiej wybrać kontener.
    \item \textbf{Zwracanie adresu lokalnej zmiennej.}
    Jej czas życia kończy się przy wyjściu z zakresu.
    \item \textbf{Przekazywanie surowego wskaźnika z \texttt{get()}
    do nowego \texttt{shared\_ptr}.} Tworzy niezależny blok kontrolny.
    \item \textbf{Nieuzasadnione \texttt{shared\_ptr} wszędzie.}
    Współdzielenie własności kosztuje i zaciera odpowiedzialność.
    \item \textbf{Cykl silnych właścicieli.} Przerywaj go przez
    \texttt{weak\_ptr} lub przeprojektowanie relacji.
    \item \textbf{Nadużywanie \texttt{release()}.} Po wywołaniu
    ktoś musi jawnie przejąć odpowiedzialność za zasób.
    \item \textbf{Brak wirtualnego destruktora w bazie.}
    Usuwanie obiektu pochodnego przez wskaźnik bazowy z domyślnym
    deleterem może wówczas prowadzić do UB.
    \item \textbf{Założenie, że \texttt{shared\_ptr} chroni obiekt
    przed wyścigami.} Chroni zarządzanie własnością, nie samą mutację.
    \item \textbf{Przechowywanie wskaźników do elementów \texttt{vector}
    po realokacji.} Przy operacjach zmieniających bufor pobieraj
    adresy ponownie.
\end{enumerate}

\section{Jak wybrać właściwy typ?}

\begin{center}
\small
\begin{tabular}{|p{0.43\textwidth}|p{0.43\textwidth}|}
\hline
\textbf{Sytuacja} & \textbf{Najczęstszy wybór} \\
\hline
Pojedynczy obiekt o prostym czasie życia & Zwykła zmienna \texttt{T} \\
\hline
Dynamiczna lista elementów & \texttt{std::vector<T>} \\
\hline
Obiekt wymagany tylko do odczytu & \texttt{const T\&} \\
\hline
Obiekt wymagany i modyfikowany & \texttt{T\&} \\
\hline
Obiekt opcjonalny bez przekazania własności & \texttt{T*} lub \texttt{const T*} \\
\hline
Pojedynczy właściciel dynamicznego obiektu & \texttt{std::unique\_ptr<T>} \\
\hline
Kilku niezależnych właścicieli & \texttt{std::shared\_ptr<T>} \\
\hline
Obserwacja obiektu współdzielonego & \texttt{std::weak\_ptr<T>} \\
\hline
Fragment istniejącej tablicy, bez własności & \texttt{std::span<T>} \\
\hline
Polimorficzna kolekcja obiektów & \texttt{std::vector} zawierający \texttt{unique\_ptr} \\
\hline
\end{tabular}
\end{center}

\subsection{Praktyczny przykład: kolekcja obiektów polimorficznych}

\begin{lstlisting}
#include <iostream>
#include <memory>
#include <vector>

struct Task {
    virtual ~Task() = default;
    virtual void run() const = 0;
};

struct PrintTask : Task {
    void run() const override { std::cout << "hello\n"; }
};

int main() {
    std::vector<std::unique_ptr<Task>> tasks;
    tasks.push_back(std::make_unique<PrintTask>());

    for (const auto& task : tasks) {
        task->run();
    }
} // vector and tasks destroy everything automatically
\end{lstlisting}

\section{Podsumowanie i praktyczne reguły}

\begin{itemize}
    \item Preferuj \emph{przechowywanie przez wartość}, jeśli
    obiekt nie wymaga odrębnej, dynamicznej tożsamości.
    \item Gdy obiekt ma jednego właściciela, używaj
    \texttt{std::unique\_ptr}; własność przenoś przez
    \texttt{std::move}.
    \item Sięgaj po \texttt{std::shared\_ptr} dopiero wtedy,
    gdy kilka części programu \emph{rzeczywiście} musi niezależnie
    utrzymywać ten sam obiekt przy życiu.
    \item Jeśli obserwator nie powinien utrzymywać obiektu przy życiu,
    użyj \texttt{std::weak\_ptr} i metody \texttt{lock()}.
    \item Do niewłasnościowego dostępu używaj
    \texttt{T\&}, \texttt{const T\&}, \texttt{T*} lub
    \texttt{std::span<T>} w zależności od kontraktu.
    \item Nie dereferencjonuj wskaźnika tylko dlatego, że jest
    różny od \texttt{nullptr}: obiekt może już nie istnieć.
    \item Preferuj \texttt{std::vector} i \texttt{std::array}
    zamiast ręcznie alokowanych tablic.
    \item Zasoby systemowe obejmuj RAII i odpowiednim deleterem;
    osobno analizuj bezpieczeństwo wątkowe wskazywanego obiektu.
\end{itemize}

\paragraph{Konkluzja.}
Najważniejsze pytanie brzmi nie ``jak zdobyć adres?'', ale
``\emph{kto jest właścicielem obiektu, jak długo on istnieje
i kto ma prawo z niego korzystać?}''. Właśnie te trzy kwestie
rozstrzygają, czy w danym miejscu potrzebujemy referencji,
surowego wskaźnika, \texttt{unique\_ptr}, \texttt{shared\_ptr}
czy \texttt{weak\_ptr}.

% Dalsza lektura:
% https://en.cppreference.com/cpp/memory/unique_ptr
% https://en.cppreference.com/cpp/memory/shared_ptr
% https://en.cppreference.com/cpp/memory/weak_ptr
% https://en.cppreference.com/cpp/memory/enable_shared_from_this
% https://en.cppreference.com/cpp/container/span


\chapter[Wskaźniki i inteligentne wskaźniki w C++]{Wskaźniki, zarządzanie pamięcią i inteligentne wskaźniki w C++}
\label{chap:wskazniki-cpp}

\section{Po co istnieją wskaźniki?}

Wskaźnik (ang. \emph{pointer}) przechowuje wartość pozwalającą
odwołać się do obiektu albo funkcji. Typowy wskaźnik do obiektu ma
postać \texttt{T*}. W odróżnieniu od obiektu typu \texttt{T},
wskaźnik może wskazywać na inny obiekt, może być pusty, a samo
wskazywanie \emph{nie oznacza posiadania wskazywanego obiektu}.

Wskaźniki wykorzystuje się między innymi do:
\begin{itemize}
    \item przekazywania i modyfikowania istniejących obiektów bez ich kopiowania;
    \item opisywania opcjonalnej relacji ``obiekt istnieje albo go nie ma'';
    \item implementowania dynamicznych struktur danych, takich jak listy,
    drzewa i grafy;
    \item realizowania polimorfizmu i dynamicznego wiązania;
    \item współpracy z interfejsami C, sterownikami i bibliotekami systemowymi;
    \item kontrolowania czasu życia zasobów za pomocą
    \texttt{std::unique\_ptr}, \texttt{std::shared\_ptr} i
    \texttt{std::weak\_ptr}.
\end{itemize}

\paragraph{Najważniejsza zasada.}
W nowoczesnym C++ wskaźnik \texttt{T*} traktujemy przede wszystkim
jako \emph{nieposiadający dostęp} do obiektu (\emph{observer}).
Odpowiedzialność za czas życia obiektu powinna być reprezentowana
osobno i jednoznacznie. Najczęściej wystarcza obiekt przechowywany
bezpośrednio, kontener standardowy albo \texttt{std::unique\_ptr}.

\section{Podstawy: adres, dereferencja i \texttt{nullptr}}

\begin{center}\small
\begin{tabularx}{\textwidth}{@{}p{0.32\textwidth}X@{}}
\toprule
Zapis & Za co odpowiada \\
\midrule
\texttt{\&x} & Pobiera adres obiektu \code{x}. \\
\code{*p} & Uzyskuje obiekt wskazywany przez \code{p}; wymaga ważnego, niepustego wskaźnika. \\
\code{p->field} & Udostępnia składową przez wskaźnik; odpowiada \code{(*p).field}. \\
\code{nullptr} & Oznacza brak wskazywanego obiektu. \\
\code{std::span<T>} & Łączy dostęp do kolejnych elementów z ich liczbą, ale nie posiada pamięci. \\
\code{std::addressof(x)} & Pobiera rzeczywisty adres obiektu nawet przy przeciążonym operatorze \texttt{\&}. \\
\bottomrule
\end{tabularx}
\end{center}

\subsection{Deklarowanie wskaźników}

Operator \texttt{\&} zastosowany do obiektu daje jego adres, natomiast
operator jednoargumentowy \texttt{*} służy do dereferencji (dostępu do
wskazywanego obiektu):

\begin{lstlisting}
#include <iostream>

int main() {
    int x = 10;
    int* p = &x;             // p points to x

    std::cout << x << '\n';  // 10
    std::cout << *p << '\n'; // 10

    *p = 25;                 // modifies x
    std::cout << x << '\n';  // 25
}
\end{lstlisting}

W deklaracji \texttt{int* p} gwiazdka należy do deklaratora, nie jest
operacją dereferencji. Nie należy też utożsamiać wskaźnika z samym
obiektem: \texttt{p} jest osobnym obiektem o własnej wartości i adresie.

\begin{lstlisting}
int x = 7;
int* p = &x;
int** pp = &p;  // pointer to a pointer

**pp = 9;       // changes x
\end{lstlisting}

Zapis \texttt{int* a, b;} deklaruje wskaźnik \texttt{a} oraz
zwykłą zmienną \texttt{b}. Dla czytelności lepiej deklarować je
oddzielnie.

\subsection{Wskaźnik pusty i poprawna dereferencja}

Pusty wskaźnik zapisujemy za pomocą \texttt{nullptr} (od C++11).
Nie wolno go dereferencjonować. Samo porównanie z
\texttt{nullptr} jest bezpieczne:

\begin{lstlisting}
int* p = nullptr;

if (p != nullptr) {
    *p = 5;
}

if (p) {                 // equivalent non-null test
    *p = 6;
}
\end{lstlisting}

\texttt{nullptr} ma osobny typ \texttt{std::nullptr\_t}; w odróżnieniu od
historycznych zapisów \texttt{0} i \texttt{NULL} nie jest zwykłą
liczbą całkowitą. Pomaga to przy przeciążaniu funkcji:

\begin{lstlisting}
void f(int);
void f(int*);

f(0);         // calls f(int)
f(nullptr);   // calls f(int*)
\end{lstlisting}

\paragraph{Uwaga o wskaźnikach niezainicjalizowanych.}
Lokalne \texttt{int* p;} bez inicjalizacji ma nieokreśloną wartość;
korzystanie z niej może prowadzić do niezdefiniowanego zachowania.
Zawsze inicjalizuj wskaźnik poprawnym adresem lub \texttt{nullptr}.

\subsection{Dostęp do składowych: \texttt{.} i \texttt{->}}

Operator \texttt{.} wybiera składową obiektu, natomiast
\texttt{->} wybiera składową przez wskaźnik:

\begin{lstlisting}
#include <string>

struct User {
    std::string name;
    int score = 0;
};

User u{"Anna", 12};
User* p = &u;

u.score = 13;
p->score = 14;       // same as (*p).score = 14
\end{lstlisting}

Nawiasy w \texttt{(*p).score} są potrzebne ze względu na priorytety
operatorów. Zapis \texttt{*p.score} nie oznacza tego samego.

\section{Wskaźniki, referencje i kwalifikator \texttt{const}}

\subsection{Kiedy wskaźnik, a kiedy referencja?}

Referencja \texttt{T\&} jest zwykle odpowiednia, gdy obiekt musi
istnieć i parametr nie jest opcjonalny. Wskaźnik \texttt{T*} nadaje
się, gdy brak obiektu jest dopuszczalny lub gdy chcemy zmieniać
obiekt, na który wskazujemy.

\begin{lstlisting}
void increase(int& x) { ++x; }

void increase_if_present(int* x) {
    if (x) { ++*x; }
}

int n = 3;
increase(n);             // n == 4
increase_if_present(&n); // n == 5
increase_if_present(nullptr); // no operation
\end{lstlisting}

Referencja również może stać się \emph{wisząca}, gdy obiekt,
do którego się odnosi, przestanie istnieć. Typ \texttt{T\&}
nie gwarantuje samodzielnie, że czas życia obiektu jest wystarczająco długi.

\subsection{Co jest stałe: wskaźnik czy obiekt?}

\begin{lstlisting}
int a = 1;
int b = 2;

const int* p = &a;       // read-only access through p
p = &b;                 // OK: p can point elsewhere
// *p = 3;              // error

int* const q = &a;       // q cannot be reassigned
*q = 3;                 // OK: a is mutable
// q = &b;              // error

const int* const r = &a; // neither reassignment nor write
int const* s = &a;       // same type as const int*
\end{lstlisting}

Reguła praktyczna: \texttt{const} po \texttt{*} dotyczy samego
wskaźnika, a przed \texttt{*} --- wskazywanego typu.

\paragraph{Pozorna niemutowalność.}
Wskaźnik \texttt{const T*} nie pozwala modyfikować obiektu
\emph{przez ten wskaźnik}. Nie oznacza, że obiekt nie może ulec zmianie
inną drogą:

\begin{lstlisting}
int x = 10;
const int* read_only = &x;
x = 20;                 // allowed
// *read_only = 30;     // not allowed
\end{lstlisting}

\section{Czas życia obiektu i wskaźniki wiszące}

\subsection{Adres nie oznacza prawa dostępu}

Wskaźnik może przechowywać adres miejsca, w którym obiekt już
nie istnieje. Jest wtedy wskaźnikiem \emph{wiszącym}
(ang. \emph{dangling pointer}). Dereferencja takiego wskaźnika
prowadzi do niezdefiniowanego zachowania (\emph{undefined behavior}, UB).

\begin{lstlisting}
int* bad() {
    int local = 42;
    return &local;       // WRONG: local dies at return
}

// int x = *bad();       // UB: dangling pointer
\end{lstlisting}

Prawidłowo można zwrócić samą wartość:

\begin{lstlisting}
int good() {
    int local = 42;
    return local;        // safe value return
}
\end{lstlisting}

\subsection{Unieważnianie wskaźników przez kontenery}

Nie każdy wskaźnik staje się wiszący przez wyjście z funkcji.
Relokacja elementów \texttt{std::vector} także może unieważnić
wskaźniki, referencje i iteratory do tych elementów:

\begin{lstlisting}
#include <vector>

std::vector<int> values{10, 20};
int* p = &values[0];

values.reserve(100);     // may reallocate storage
// *p = 50;              // unsafe if reallocation occurred
p = &values[0];          // reacquire a valid pointer
\end{lstlisting}

Po zmianie pojemności wymagającej realokacji adresy elementów
się zmieniają. W ogólności po operacji mogącej unieważnić wskaźnik
należy pobrać go ponownie, a nie próbować odgadnąć, czy nadal działa.

\subsection{Trzy niezależne pytania}

Przy każdym wskaźniku rozstrzygaj:
\begin{enumerate}
    \item Czy wartość wskaźnika jest poprawna (np. nie jest
    niezainicjalizowana)?
    \item Czy obiekt, na który wskazujemy, nadal istnieje?
    \item Czy operacja jest dozwolona dla danego typu i zakresu pamięci
    (np. czy nie wychodzi poza tablicę)?
\end{enumerate}

Sam test \texttt{p != nullptr} rozwiązuje co najwyżej fragment
pierwszego problemu, ale \emph{nie} wykrywa wiszących wskaźników.

\section{Tablice, arytmetyka wskaźników i \texttt{std::span}}

\subsection{Arytmetyka wskaźników}

Dla wskaźnika do elementu tablicy operacje \texttt{p + k} i
\texttt{p[k]} odnoszą się do jej elementów, pod warunkiem że
poruszamy się w dozwolonym zakresie. Można utworzyć wskaźnik
\emph{jeden za końcem} tablicy, ale nie wolno go dereferencjonować.

\begin{lstlisting}
int a[]{4, 8, 12};
int* p = a;          // points to a[0]

int x = *(p + 1);    // 8
int y = p[2];       // 12
int* end = a + 3;   // valid one-past-the-end pointer
// int z = *end;     // UB
\end{lstlisting}

Dodawanie i odejmowanie wskaźników ma sens w obrębie \emph{tej samej}
tablicy (z uwzględnieniem pozycji jeden za końcem). Nie należy
wykorzystywać takich działań do przemieszczania się między dowolnymi
obiektami w pamięci.

\subsection{Tablica to nie wskaźnik}

Typ \texttt{int[3]} i typ \texttt{int*} są różne, choć w wielu
wyrażeniach tablica jest automatycznie konwertowana na wskaźnik
do pierwszego elementu (tzw. \emph{array-to-pointer decay}).

\begin{lstlisting}
int a[3] = {1, 2, 3};
int* p = a;

static_assert(sizeof(a) == 3 * sizeof(int));
// sizeof(p) is the size of a pointer, not of the array
\end{lstlisting}

Funkcja \texttt{void f(int* p)} nie zna długości tablicy.
Jeśli trzeba pracować z ciągłym fragmentem pamięci o znanym
rozmiarze, C++20 udostępnia \texttt{std::span}:

\begin{lstlisting}
#include <span>
#include <vector>

int sum(std::span<const int> xs) {
    int result = 0;
    for (int x : xs) {
        result += x;
    }
    return result;
}

int a[]{1, 2, 3};
std::vector<int> v{4, 5, 6};

int s1 = sum(a);  // 6
int s2 = sum(v);  // 15
\end{lstlisting}

\texttt{std::span<T>} przechowuje widok ciągłych danych
(z informacją o liczbie elementów), ale \emph{nie jest właścicielem}
pamięci. Gdy dane znikną albo przeniosą się w inne miejsce,
\texttt{span} również może stać się niepoprawny. Dla większości
dynamicznych tablic własnościowych preferuj \texttt{std::vector<T>}.

\section{\texttt{void*}, konwersje i rzutowania wskaźników}

\subsection{\texttt{void*}: adres bez informacji o typie}

Wskaźnik \texttt{void*} może przechowywać adres obiektu bez
informacji o jego konkretnym typie. Bywa używany w interfejsach C
i bibliotekach niskopoziomowych, ale nie pozwala na bezpośrednią
dereferencję:

\begin{lstlisting}
int x = 42;
void* opaque = &x;              // type erased
// *opaque = 7;                 // error: cannot dereference void*

int* p = static_cast<int*>(opaque);
*p = 7;                        // valid: original object is int
\end{lstlisting}

Przywrócenie błędnego typu i późniejsza próba użycia obiektu
mogą prowadzić do niezdefiniowanego zachowania. Jeśli typ jest
znany, zwykle lepiej przechowywać od razu \texttt{T*}.

\subsection{\texttt{static\_cast}, \texttt{dynamic\_cast} i
\texttt{reinterpret\_cast}}

\begin{itemize}
    \item \texttt{static\_cast<T*>} nadaje się m.in. do poprawnych
    konwersji pomiędzy \texttt{void*} a wskaźnikiem obiektowym,
    jeśli rzeczywisty typ i wyrównanie na to pozwalają.
    \item \texttt{dynamic\_cast<T*>} służy do sprawdzanych w czasie
    wykonania konwersji w hierarchiach klas polimorficznych.
    Niepowodzenie konwersji wskaźnika zwraca \texttt{nullptr}.
    \item \texttt{reinterpret\_cast<T*>} umożliwia niskopoziomową
    reinterpretację reprezentacji wskaźnika, ale sama konwersja
    nie uprawnia do dowolnej dereferencji ani nie tworzy obiektu
    odpowiedniego typu.
\end{itemize}

\begin{lstlisting}
struct Base {
    virtual ~Base() = default;
};

struct Derived : Base {
    int value = 12;
};

Derived d;
Base* base = &d;

if (Derived* p = dynamic_cast<Derived*>(base)) {
    p->value = 30;          // checked downcast succeeded
}
\end{lstlisting}

Przy przekazywaniu parametrów funkcji nie używaj rzutowań
jako sposobu na obejście \texttt{const}, sprawdzania czasu życia
ani zasad własności. Dobrze zaprojektowany interfejs zazwyczaj
w ogóle nie wymaga rzutowania wskaźników.

\section{Pamięć dynamiczna: \texttt{new}, \texttt{delete} i ich ograniczenia}

\subsection{Klasyczny mechanizm}

Operator \texttt{new} tworzy obiekt w pamięci przydzielonej
dynamicznie i zwraca wskaźnik. Odpowiadający mu
\texttt{delete} niszczy obiekt i zwalnia pamięć:

\begin{lstlisting}
int* p = new int(42);
// use *p

delete p;
p = nullptr;           // avoids this one stale alias
\end{lstlisting}

Dla tablic obowiązuje para \texttt{new[]} i \texttt{delete[]}:

\begin{lstlisting}
int* a = new int[5]{}; // zero-initialized elements
// use a[0], ..., a[4]

delete[] a;
\end{lstlisting}

\textbf{Nie wolno} łączyć \texttt{new[]} z \texttt{delete} ani
zwykłego \texttt{new} z \texttt{delete[]}. Ponowne
\texttt{delete} tego samego żywego przydziału również jest błędem.
\texttt{delete nullptr;} jest natomiast bezpieczne.

\subsection{Dlaczego ręczne zarządzanie pamięcią jest niebezpieczne?}

\begin{lstlisting}
void fragile() {
    int* p = new int(42);
    operation_that_may_throw();
    delete p;   // skipped if the operation throws
}
\end{lstlisting}

Jeśli funkcja zgłosi wyjątek, pamięć zostaje utracona
(\emph{memory leak}). Analogiczny problem powstaje przy
wielu instrukcjach \texttt{return}, kopiowaniu wskaźników i
niespójnym przekazywaniu odpowiedzialności za \texttt{delete}.

\paragraph{Zalecenie.}
W typowym kodzie aplikacyjnym nie pisz bezpośrednio
\texttt{new}/\texttt{delete}. Użyj obiektu lokalnego,
\texttt{std::vector<T>} albo inteligentnego wskaźnika.
Ręczne operacje alokowania pozostają potrzebne w implementacjach
niskopoziomowych, alokatorach i przy integracji z określonym API.

\section{RAII i semantyka własności}

\subsection{RAII: zasób należy do obiektu}

\emph{Resource Acquisition Is Initialization} (RAII) oznacza
powiązanie przejęcia zasobu z utworzeniem obiektu zarządzającego,
a zwolnienia zasobu z jego destruktorem. Stosuje się je nie tylko
do pamięci, ale też do plików, blokad, gniazd sieciowych i uchwytów
systemowych.

\begin{lstlisting}
#include <fstream>
#include <string>

void save(const std::string& path) {
    std::ofstream file(path);
    file << "example\n";
} // file is closed by its destructor
\end{lstlisting}

Destruktor uruchamia się podczas normalnego opuszczania zakresu
oraz podczas odwijania stosu po wyjątku. Stąd RAII upraszcza
bezpieczeństwo wyjątkowe.

\subsection{Trzy rodzaje relacji z obiektem}

\begin{description}
    \item[Własność wyłączna (\emph{unique ownership}).]
    Dokładnie jeden obiekt zarządzający odpowiada za zniszczenie
    zasobu: najczęściej \texttt{std::unique\_ptr<T>}.
    \item[Własność współdzielona (\emph{shared ownership}).]
    Wielu właścicieli może utrzymywać jeden zasób przy życiu:
    \texttt{std::shared\_ptr<T>}.
    \item[Dostęp bez własności (\emph{non-owning access}).]
    Użytkownik widzi obiekt, ale nie przedłuża jego życia:
    zwykły wskaźnik, referencja, \texttt{std::span<T>} lub
    \texttt{std::weak\_ptr<T>} (ten ostatni obserwuje stan
    zarządzany przez \texttt{shared\_ptr} i pozwala warunkowo
    uzyskać właściciela).
\end{description}

Dobór reprezentacji wynika z \emph{kontraktu własności}, a nie
z samego faktu, że obiekt znajduje się w pamięci dynamicznej.

\section{\texttt{std::unique\_ptr}: jeden właściciel}

\begin{center}\small
\begin{tabularx}{\textwidth}{@{}p{0.36\textwidth}X@{}}
\toprule
Mechanizm & Własność i czas życia \\
\midrule
\code{std::make_unique<T>(...)} & Tworzy obiekt i zwraca jego jedynego właściciela. \\
\code{unique_ptr::get()} & Zwraca surowy wskaźnik bez przekazania własności. \\
\code{unique_ptr::reset(p)} & Niszczy obecnie posiadany obiekt i przejmuje wskaźnik \code{p}. \\
\code{unique_ptr::release()} & Oddaje surowy wskaźnik bez zniszczenia obiektu; odtąd trzeba osobno zadbać o jego zwolnienie. \\
\code{std::shared_ptr<T>} & Współdzieli własność; obiekt znika po odejściu ostatniego właściciela. \\
\code{weak_ptr::lock()} & Próbuje uzyskać \code{shared_ptr}; zwraca pusty wynik, jeśli obiekt wygasł. \\
\bottomrule
\end{tabularx}
\end{center}

\subsection{Tworzenie obiektu i automatyczne niszczenie}

\texttt{std::unique\_ptr<T>} (nagłówek \texttt{<memory>}, C++11)
zarządza pojedynczym obiektem przez wyłączną własność. Od C++14
najczęściej tworzymy go przez \texttt{std::make\_unique<T>}:

\begin{lstlisting}
#include <memory>
#include <string>

#include <utility>

struct Person {
    explicit Person(std::string n) : name(std::move(n)) {}
    std::string name;
};

void example() {
    auto p = std::make_unique<Person>("Ala");
    p->name = "Ola";
    // no delete needed
} // destructor destroys Person
\end{lstlisting}

\texttt{unique\_ptr} można traktować jako obiekt zarządzający:
jego destruktor wywołuje odpowiednią funkcję usuwającą zasób
(\emph{deleter}). Domyślnie jest nią \texttt{std::default\_delete<T>}.

\subsection{Brak kopiowania i przenoszenie własności}

Wskaźnik o własności wyłącznej nie może być kopiowany.
Może natomiast zostać \emph{przeniesiony} za pomocą
\texttt{std::move}:

\begin{lstlisting}
#include <memory>
#include <utility>

auto first = std::make_unique<int>(10);
// auto bad = first;                  // error: copy

auto second = std::move(first);      // transfers ownership
// first == nullptr
*second = 20;
\end{lstlisting}

\texttt{std::move} samo niczego nie zwalnia i nie przenosi fizycznie
wskazywanego obiektu. Zwraca wyrażenie pozwalające wywołać konstruktor
lub przypisanie przenoszące. W tym przypadku przenoszona jest
\emph{własność wskaźnika}, a obiekt typu \texttt{int} nie zmienia adresu.

\subsection{Przekazywanie do funkcji: kto jest właścicielem?}

Rozróżnijmy trzy interfejsy:

\begin{lstlisting}
#include <memory>
#include <utility>

struct Widget { void run() {} };

void inspect(const Widget& w);             // borrows, non-null
void maybe_inspect(const Widget* w);       // borrows, nullable
void take(std::unique_ptr<Widget> w);      // takes ownership

void demo() {
    auto p = std::make_unique<Widget>();
    inspect(*p);                // p still owns the object
    maybe_inspect(p.get());     // p still owns the object
    take(std::move(p));         // ownership transferred
    // p is now empty
}
\end{lstlisting}

Jeżeli funkcja \emph{jedynie używa} obiektu, przekazywanie przez
\texttt{T\&}, \texttt{const T\&} lub \texttt{T*} jest zwykle
czytelniejsze niż przekazywanie \texttt{unique\_ptr} przez referencję.
Parametr \texttt{unique\_ptr<T>} przyjmowany \emph{przez wartość}
sygnalizuje przejęcie własności. Dla funkcji, która ma zastąpić
obiekt będący własnością wywołującego, można użyć
\texttt{std::unique\_ptr<T>\&}.

\subsection{Metody \texttt{get}, \texttt{reset}, \texttt{release}, \texttt{swap}}

\begin{lstlisting}
#include <memory>
#include <utility>

auto p = std::make_unique<int>(7);
int* observer = p.get();       // borrow; do not delete

p.reset(new int(9));           // destroys old int
// observer is now dangling

auto q = std::make_unique<int>(11);
p.swap(q);                    // exchanges managed objects

int* raw = q.release();        // q becomes empty; raw owns now
// responsibility for delete raw is now explicit
delete raw;
\end{lstlisting}

\texttt{get()} zwraca dostęp bez własności. \texttt{reset()} zastępuje
zarządzany obiekt, zwalniając dotychczasowy. \texttt{release()} oddaje
surowy wskaźnik \emph{bez zwolnienia obiektu}; dlatego jest potrzebne
głównie przy przekazywaniu własności do interfejsu, który tego wymaga.
Nie należy mylić \texttt{release()} z \texttt{reset()}.

\subsection{Tablice: \texttt{unique\_ptr<T[]>} a \texttt{vector<T>}}

\begin{lstlisting}
#include <memory>
#include <vector>

auto a = std::make_unique<int[]>(5); // array of 5 zeros
a[2] = 17;

std::vector<int> b(5);         // usually more convenient
b.at(2) = 17;                  // checked indexing
\end{lstlisting}

\texttt{unique\_ptr<T[]>} dobiera \texttt{delete[]} i udostępnia
\texttt{operator[]}, lecz sam nie przechowuje rozmiaru.
Zwykle preferuj \texttt{std::vector<T>}, bo posiada rozmiar,
iteratory i bogatszy interfejs. Dla tablicy o stałym rozmiarze
określonym w kodzie rozważ \texttt{std::array<T, N>}.

\subsection{Własny deleter: zasób inny niż pamięć \texttt{new}}

Inteligentny wskaźnik może wywoływać dowolną funkcję zwalniającą
zasób. Typowym przykładem jest plik w interfejsie C:

\begin{lstlisting}
#include <cstdio>
#include <memory>

struct FileCloser {
    void operator()(std::FILE* file) const noexcept {
        if (file) { std::fclose(file); }
    }
};

using FilePtr = std::unique_ptr<std::FILE, FileCloser>;

void example() {
    FilePtr file(std::fopen("log.txt", "w"));
    if (!file) { return; }

    std::fputs("hello\n", file.get());
} // fclose is called automatically
\end{lstlisting}

W praktyce deleter bywa osobną strukturą, a jego typ wchodzi
w skład typu \texttt{unique\_ptr<T, Deleter>}. Mechanizm pozwala
owinąć RAII również zasoby zwalniane przez \texttt{free()},
\texttt{fclose()} lub bibliotekę zewnętrzną.

\subsection{Polimorfizm i wirtualny destruktor}

\begin{lstlisting}
#include <memory>

struct Shape {
    virtual ~Shape() = default;
    virtual double area() const = 0;
};

struct Circle : Shape {
    explicit Circle(double r) : radius(r) {}
    double area() const override { return 3.14159 * radius * radius; }
    double radius;
};

std::unique_ptr<Shape> make_shape() {
    return std::make_unique<Circle>(2.0);
}
\end{lstlisting}

Przy niszczeniu obiektu pochodnego przez wskaźnik do bazy
zwykle potrzebny jest wirtualny destruktor bazy.
Dotyczy to również przenoszenia własności z klasy pochodnej
do bazowej przy użyciu \texttt{unique\_ptr} i domyślnego deletera.

\section{\texttt{std::shared\_ptr}: współdzielona własność}

\subsection{Jak działa licznik właścicieli?}

\texttt{std::shared\_ptr<T>} (C++11) udostępnia współdzieloną
własność zasobu. Zwykle tworzymy go za pomocą
\texttt{std::make\_shared<T>}:

\begin{lstlisting}
#include <memory>

auto a = std::make_shared<int>(10);
auto b = a;             // shared ownership

// a.use_count() == 2 (provided there are no other owners)
b.reset();              // releases b's ownership
// a.use_count() == 1
\end{lstlisting}

Gdy liczba \emph{silnych właścicieli} spadnie do zera, zarządzany
obiekt zostaje zniszczony. Implementacja utrzymuje tzw.
\emph{blok kontrolny} (\emph{control block}) zawierający między
innymi liczniki właścicieli, informacje potrzebne do zniszczenia
zasobu i ewentualny deleter.

\paragraph{Schemat ideowy.}
\begin{lstlisting}
shared_ptr a ---\
                +--> [ control block: strong = 2 ] --> [ T ]
shared_ptr b ---/           weak = 0
\end{lstlisting}

\texttt{use\_count()} ma znaczenie przede wszystkim diagnostyczne.
Nie powinno się na podstawie odczytanego licznika podejmować decyzji
o bezpieczeństwie wątkowym ani o wyłącznej własności, która może
zmienić się równolegle.

\subsection{Tworzenie i przekazywanie \texttt{shared\_ptr}}

\begin{lstlisting}
#include <memory>

struct Data { int x = 0; };

std::shared_ptr<Data> create() {
    return std::make_shared<Data>();
}

void keep(std::shared_ptr<Data> p) {
    // stores a copy of p somewhere, if needed
}

void read_only(const Data& data) {
    // does not require shared ownership
}

void demo() {
    auto p = create();
    read_only(*p); // borrow
    keep(p);       // creates another temporary shared owner
}
\end{lstlisting}

Parametr \texttt{shared\_ptr<T>} przyjmowany przez wartość oznacza,
że funkcja otrzymuje własny udział we własności na czas wywołania
i ewentualnie może go przechować. Jeśli funkcja wyłącznie
czyta obiekt, zwykle wystarcza \texttt{const T\&} lub \texttt{T*}.
Samo przekazanie \texttt{const shared\_ptr<T>\&} unika inkrementacji
licznika, ale bardziej wiąże interfejs z modelem własności.

\subsection{\texttt{make\_shared} a konstruktor ze wskaźnika}

W typowych implementacjach \texttt{make\_shared} umieszcza obiekt
i blok kontrolny w jednym przydziale pamięci, co może ograniczać
narzut alokacji. Natomiast konstrukcja \texttt{shared\_ptr<T>(new T)}
zwykle wymaga osobnych przydziałów.

\begin{lstlisting}
auto preferred = std::make_shared<int>(42);
std::shared_ptr<int> also_valid(new int(42));
\end{lstlisting}

Nie wolno jednak tworzyć dwóch \texttt{shared\_ptr} niezależnie
z \emph{tego samego} surowego wskaźnika:

\begin{lstlisting}
int* raw = new int(42);
std::shared_ptr<int> a(raw);
// std::shared_ptr<int> b(raw); // WRONG: second control block
\end{lstlisting}

Dwa niezależne bloki kontrolne próbowałyby zwolnić ten sam obiekt.
Poprawne współdzielenie powstaje przez \emph{kopiowanie istniejącego}
\texttt{shared\_ptr}: \texttt{auto b = a;}.

\paragraph{Kiedy nie stosować \texttt{make\_shared}?}
Zintegrowany przydział może zatrzymywać pamięć zajmowaną przez
obiekt w przydzielonym bloku do momentu zniszczenia także ostatniego
\texttt{weak\_ptr}, nawet gdy destruktor obiektu już się wykonał.
Przy dużych obiektach i długo żyjących słabych referencjach
warto rozważyć osobny przydział. Innymi powodami mogą być
specjalne wymagania alokacji lub konstruktory niedostępne
dla \texttt{make\_shared}.

\subsection{Współdzielenie własności a współdzielenie mutacji}

Sam \texttt{shared\_ptr} nie synchronizuje dostępu do pól obiektu:

\begin{lstlisting}
#include <memory>
#include <mutex>

struct Counter {
    std::mutex mutex;
    int value = 0;

    void increment() {
        std::lock_guard<std::mutex> lock(mutex);
        ++value;
    }
};

auto p = std::make_shared<Counter>();
// many threads may own copies of p, but must synchronize
// concurrent modifications of Counter::value
\end{lstlisting}

Blok kontrolny pozwala różnym egzemplarzom \texttt{shared\_ptr}
bezpiecznie współdzielić własność. Równoczesna modyfikacja
\emph{tego samego} egzemplarza wymaga jednak synchronizacji.
Od C++20 można użyć \texttt{std::atomic<std::shared\_ptr<T>>}.

\section{\texttt{std::weak\_ptr}: obserwacja bez utrzymywania przy życiu}

\subsection{Problem cykli własności}

Gdy obiekty wzajemnie posiadają się przez \texttt{shared\_ptr},
liczniki właścicieli mogą nigdy nie spaść do zera:

\begin{lstlisting}
#include <memory>

struct Node {
    std::shared_ptr<Node> next;
};

void leak_cycle() {
    auto a = std::make_shared<Node>();
    auto b = std::make_shared<Node>();
    a->next = b;
    b->next = a;          // cycle of strong ownership
} // nodes remain alive after both locals are destroyed
\end{lstlisting}

To nie jest problem \emph{niezerowego wskaźnika}, lecz błędnego
modelu własności: każdy węzeł utrzymuje przy życiu drugi.

\subsection{Przerywanie cyklu i \texttt{lock()}}

\texttt{std::weak\_ptr<T>} obserwuje ten sam blok kontrolny,
ale nie zwiększa liczby silnych właścicieli:

\begin{lstlisting}
#include <memory>

struct Node {
    std::weak_ptr<Node> parent;         // non-owning
    std::shared_ptr<Node> child;        // owning
};

void example() {
    auto root = std::make_shared<Node>();
    auto leaf = std::make_shared<Node>();

    root->child = leaf;
    leaf->parent = root;

    if (auto parent = leaf->parent.lock()) {
        // parent is a temporary shared owner
        (void)parent;
    }
} // no ownership cycle
\end{lstlisting}

Metoda \texttt{lock()} atomowo próbuje uzyskać
\texttt{shared\_ptr<T>}. Jeżeli obiekt jeszcze istnieje,
zwraca właściciela utrzymującego go przy życiu przez czas
używania. W przeciwnym razie zwraca pusty \texttt{shared\_ptr}.

\begin{lstlisting}
std::weak_ptr<int> observer;
{
    auto p = std::make_shared<int>(8);
    observer = p;
    if (auto alive = observer.lock()) {
        *alive += 1;           // safe while alive owns it
    }
}

// observer.expired() is now true
// observer.lock() now returns an empty shared_ptr
\end{lstlisting}

W kodzie współbieżnym nie stosuj schematu
\texttt{if (!w.expired()) { ... }} z późniejszym odwołaniem:
między sprawdzeniem a użyciem obiekt może zniknąć.
Preferuj jedną operację \texttt{lock()} i posługuj się
otrzymanym \texttt{shared\_ptr}.

\section{\texttt{enable\_shared\_from\_this} i wskaźnik \texttt{this}}

Gdy metoda obiektu zarządzanego przez \texttt{shared\_ptr}
chce przekazać dalej \emph{współdzieloną własność do siebie},
nie może wykonać \texttt{std::shared\_ptr<T>(this)}:
powstałby drugi, niezależny blok kontrolny.

Można zamiast tego dziedziczyć publicznie po
\texttt{std::enable\_shared\_from\_this<T>}:

\begin{lstlisting}
#include <memory>

struct Session : std::enable_shared_from_this<Session> {
    std::shared_ptr<Session> self() {
        return shared_from_this();
    }
};

void example() {
    auto session = std::make_shared<Session>();
    auto second_owner = session->self();
    // both share the same control block
}
\end{lstlisting}

Metodę \texttt{shared\_from\_this()} wolno wywołać dopiero wtedy,
gdy obiektem zarządza już \texttt{shared\_ptr}. Jej wywołanie
na obiekcie lokalnym albo w konstruktorze, zanim powstanie
właściciel, zgłasza wyjątek \texttt{std::bad\_weak\_ptr}.

Od C++17 dostępna jest również metoda \texttt{weak\_from\_this()}.

\section{Zaawansowany przykład: wskaźnik aliasujący \texttt{shared\_ptr}}

\texttt{shared\_ptr} rozróżnia \emph{obiekt zarządzany} (którego
czas życia kontroluje blok kontrolny) i \emph{przechowywany
wskaźnik} (zwracany przez \texttt{get()}). Specjalny konstruktor
aliasujący pozwala udostępnić adres składowej, jednocześnie
utrzymując przy życiu cały obiekt:

\begin{lstlisting}
#include <memory>

struct Pair {
    int first = 1;
    int second = 2;
};

void example() {
    auto whole = std::make_shared<Pair>();
    std::shared_ptr<int> member(whole, &whole->second);

    whole.reset();       // Pair is still alive
    *member = 20;        // accesses Pair::second safely
} // member destroys the last strong owner
\end{lstlisting}

Konstruktor aliasujący jest przydatny np. przy udostępnianiu
podobiektów, ale wymaga prawidłowej relacji czasów życia.
Nie należy go traktować jako sposobu na naprawienie
wcześniej wiszącego wskaźnika.

\section{Wskaźniki w interfejsach funkcji i klas}

\subsection{Katalog typowych kontraktów}

\begin{center}
\small
\begin{tabular}{|p{0.31\textwidth}|p{0.56\textwidth}|}
\hline
\textbf{Interfejs} & \textbf{Znaczenie} \\
\hline
\texttt{void f(const T\&)} & Obiekt wymagany, tylko odczyt; brak przekazania własności. \\
\hline
\texttt{void f(T\&)} & Obiekt wymagany, możliwa modyfikacja; brak przekazania własności. \\
\hline
\texttt{void f(const T*)} & Obiekt opcjonalny, tylko odczyt; brak własności. \\
\hline
\texttt{void f(T*)} & Obiekt opcjonalny, możliwa modyfikacja; brak własności. \\
\hline
\texttt{void f(unique\_ptr<T>)} & Przejęcie wyłącznej własności. \\
\hline
\texttt{unique\_ptr<T> f()} & Przekazanie wyłącznej własności wywołującemu. \\
\hline
\texttt{void f(shared\_ptr<T>)} & Uzyskanie udziału we współdzielonej własności. \\
\hline
\texttt{void f(weak\_ptr<T>)} & Możliwość obserwacji obiektu, którego życie kontrolują inni. \\
\hline
\texttt{void f(span<const T>)} & Widok na ciągłą sekwencję, bez własności. \\
\hline
\end{tabular}
\end{center}

W tabeli pominięto przestrzeń nazw \texttt{std::} dla czytelności.
To \emph{konwencje projektowania API}, a nie automatyczne
gwarancje długowieczności surowych wskaźników i referencji.

\subsection{Dlaczego nie zawsze używać \texttt{shared\_ptr}?}

Rozważmy funkcję, która wyświetla nazwę:

\begin{lstlisting}
#include <iostream>
#include <memory>
#include <string>

struct Product { std::string name; };

void print(const Product& p) {
    std::cout << p.name << '\n';
}

void example() {
    Product local{"book"};
    print(local);                  // no allocation

    auto owned = std::make_unique<Product>(Product{"pen"});
    print(*owned);                 // still no extra owner
}
\end{lstlisting}

Funkcja \texttt{print} nie potrzebuje znać strategii
przechowywania obiektu. Przekazanie referencji oddziela
\emph{sposób dostępu} od \emph{sposobu zarządzania czasem życia}.

\section{PImpl: ukrywanie implementacji przez \texttt{unique\_ptr}}

Technika \emph{pointer to implementation} pozwala utrzymać
szczegóły implementacji poza publicznym nagłówkiem.
Typ implementacyjny może być niekompletny w nagłówku, ale
miejsce wykonania domyślnego deletera musi widzieć jego definicję.
Dlatego destruktor klasy definiujemy w pliku \texttt{.cpp}.

\begin{lstlisting}
// widget.hpp
#pragma once
#include <memory>

class Widget {
public:
    Widget();
    ~Widget();                 // defined in .cpp

    Widget(Widget&&) noexcept;
    Widget& operator=(Widget&&) noexcept;

    Widget(const Widget&) = delete;
    Widget& operator=(const Widget&) = delete;

    void work();

private:
    struct Impl;
    std::unique_ptr<Impl> impl_;
};
\end{lstlisting}

\begin{lstlisting}
// widget.cpp
#include "widget.hpp"
#include <memory>

struct Widget::Impl {
    int value = 0;
    void run() { ++value; }
};

Widget::Widget() : impl_(std::make_unique<Impl>()) {}
Widget::~Widget() = default;
Widget::Widget(Widget&&) noexcept = default;
Widget& Widget::operator=(Widget&&) noexcept = default;

void Widget::work() { impl_->run(); }
\end{lstlisting}

Ten przykład prezentuje semantykę przenoszenia. Po przeniesieniu
źródłowy obiekt \texttt{Widget} ma pusty \texttt{impl\_};
nie wolno wywoływać na nim \texttt{work()} bez odpowiedniego
sprawdzenia. W produkcyjnym API można ten stan jawnie obsłużyć
lub określić w kontrakcie.

\section{Wskaźniki do funkcji i składowych klas}

Wskaźniki nie muszą odnosić się do danych. Można wskazywać
funkcje, a także składowe klas, choć te ostatnie wymagają
innej składni.

\subsection{Wskaźnik do funkcji}

\begin{lstlisting}
int add(int a, int b) { return a + b; }
int mul(int a, int b) { return a * b; }

using Operation = int(*)(int, int);

int apply(int x, int y, Operation op) {
    return op(x, y);
}

int a = apply(2, 3, &add); // 5
int b = apply(2, 3, &mul); // 6
\end{lstlisting}

Obiekty wywoływalne różnych typów można przechowywać za pomocą
\texttt{std::function}. Prosty szablon lub lambda bywają jednak
bardziej efektywne.

\subsection{Wskaźnik do składowej}

\begin{lstlisting}
struct Record {
    int value = 0;
    int twice() const { return 2 * value; }
};

int Record::* field = &Record::value;
int (Record::* method)() const = &Record::twice;

Record r{7};
int x = r.*field;          // 7
int y = (r.*method)();     // 14

Record* p = &r;
int z = p->*field;         // 7
\end{lstlisting}

Wskaźnik do składowej nie jest zwykłym adresem gotowego obiektu.
Dopiero w połączeniu z konkretną instancją klasy pozwala
odwołać się do odpowiedniej składowej.

\section{Wskaźniki i wielowątkowość}

\subsection{\texttt{shared\_ptr} nie eliminuje wyścigów danych}

Możliwość bezpiecznego utrzymania czasu życia obiektu nie oznacza,
że jego pola można dowolnie zmieniać z wielu wątków.
Jeśli dwa wątki zapisują do \texttt{*p}, potrzebna jest osobna
synchronizacja (np. \texttt{std::mutex} albo atomiki).

Od C++20 dostępna jest specjalizacja
\texttt{std::atomic<std::shared\_ptr<T>>}, przydatna do atomowego
publikowania nowej wersji współdzielonego obiektu:

\begin{lstlisting}
#include <atomic>
#include <memory>

struct Config {
    int workers = 4;
};

std::atomic<std::shared_ptr<const Config>> current{
    std::make_shared<const Config>(Config{4})
};

void publish(int count) {
    current.store(std::make_shared<const Config>(Config{count}));
}

int read_workers() {
    auto snapshot = current.load(); // keeps snapshot alive
    return snapshot->workers;
}
\end{lstlisting}

Tu czytelnicy otrzymują stabilny, niemutowalny \emph{snapshot}.
Atomowość dotyczy odczytu i zapisu \texttt{shared\_ptr},
a nie arbitralnych operacji na jego zawartości.

\section{Najczęstsze błędy zebrane w jednym miejscu}

\begin{enumerate}
    \item \textbf{Dereferencja \texttt{nullptr} lub wskaźnika
    niezainicjalizowanego.} Zawsze sprawdzaj kontrakt i inicjalizuj
    wskaźniki.
    \item \textbf{Wskaźnik wiszący.} \texttt{get()}, \texttt{span}
    i surowa referencja nie przedłużają życia zasobu.
    \item \textbf{Podwójne \texttt{delete}.} Jeden przydział musi mieć
    jedną, jednoznaczną strategię zwalniania.
    \item \textbf{Mieszanie \texttt{new[]} i \texttt{delete}.}
    Do tablic służy \texttt{delete[]}; zwykle lepiej wybrać kontener.
    \item \textbf{Zwracanie adresu lokalnej zmiennej.}
    Jej czas życia kończy się przy wyjściu z zakresu.
    \item \textbf{Przekazywanie surowego wskaźnika z \texttt{get()}
    do nowego \texttt{shared\_ptr}.} Tworzy niezależny blok kontrolny.
    \item \textbf{Nieuzasadnione \texttt{shared\_ptr} wszędzie.}
    Współdzielenie własności kosztuje i zaciera odpowiedzialność.
    \item \textbf{Cykl silnych właścicieli.} Przerywaj go przez
    \texttt{weak\_ptr} lub przeprojektowanie relacji.
    \item \textbf{Nadużywanie \texttt{release()}.} Po wywołaniu
    ktoś musi jawnie przejąć odpowiedzialność za zasób.
    \item \textbf{Brak wirtualnego destruktora w bazie.}
    Usuwanie obiektu pochodnego przez wskaźnik bazowy z domyślnym
    deleterem może wówczas prowadzić do UB.
    \item \textbf{Założenie, że \texttt{shared\_ptr} chroni obiekt
    przed wyścigami.} Chroni zarządzanie własnością, nie samą mutację.
    \item \textbf{Przechowywanie wskaźników do elementów \texttt{vector}
    po realokacji.} Przy operacjach zmieniających bufor pobieraj
    adresy ponownie.
\end{enumerate}

\section{Jak wybrać właściwy typ?}

\begin{center}
\small
\begin{tabular}{|p{0.43\textwidth}|p{0.43\textwidth}|}
\hline
\textbf{Sytuacja} & \textbf{Najczęstszy wybór} \\
\hline
Pojedynczy obiekt o prostym czasie życia & Zwykła zmienna \texttt{T} \\
\hline
Dynamiczna lista elementów & \texttt{std::vector<T>} \\
\hline
Obiekt wymagany tylko do odczytu & \texttt{const T\&} \\
\hline
Obiekt wymagany i modyfikowany & \texttt{T\&} \\
\hline
Obiekt opcjonalny bez przekazania własności & \texttt{T*} lub \texttt{const T*} \\
\hline
Pojedynczy właściciel dynamicznego obiektu & \texttt{std::unique\_ptr<T>} \\
\hline
Kilku niezależnych właścicieli & \texttt{std::shared\_ptr<T>} \\
\hline
Obserwacja obiektu współdzielonego & \texttt{std::weak\_ptr<T>} \\
\hline
Fragment istniejącej tablicy, bez własności & \texttt{std::span<T>} \\
\hline
Polimorficzna kolekcja obiektów & \texttt{std::vector} zawierający \texttt{unique\_ptr} \\
\hline
\end{tabular}
\end{center}

\subsection{Praktyczny przykład: kolekcja obiektów polimorficznych}

\begin{lstlisting}
#include <iostream>
#include <memory>
#include <vector>

struct Task {
    virtual ~Task() = default;
    virtual void run() const = 0;
};

struct PrintTask : Task {
    void run() const override { std::cout << "hello\n"; }
};

int main() {
    std::vector<std::unique_ptr<Task>> tasks;
    tasks.push_back(std::make_unique<PrintTask>());

    for (const auto& task : tasks) {
        task->run();
    }
} // vector and tasks destroy everything automatically
\end{lstlisting}

\section{Podsumowanie i praktyczne reguły}

\begin{itemize}
    \item Preferuj \emph{przechowywanie przez wartość}, jeśli
    obiekt nie wymaga odrębnej, dynamicznej tożsamości.
    \item Gdy obiekt ma jednego właściciela, używaj
    \texttt{std::unique\_ptr}; własność przenoś przez
    \texttt{std::move}.
    \item Sięgaj po \texttt{std::shared\_ptr} dopiero wtedy,
    gdy kilka części programu \emph{rzeczywiście} musi niezależnie
    utrzymywać ten sam obiekt przy życiu.
    \item Jeśli obserwator nie powinien utrzymywać obiektu przy życiu,
    użyj \texttt{std::weak\_ptr} i metody \texttt{lock()}.
    \item Do niewłasnościowego dostępu używaj
    \texttt{T\&}, \texttt{const T\&}, \texttt{T*} lub
    \texttt{std::span<T>} w zależności od kontraktu.
    \item Nie dereferencjonuj wskaźnika tylko dlatego, że jest
    różny od \texttt{nullptr}: obiekt może już nie istnieć.
    \item Preferuj \texttt{std::vector} i \texttt{std::array}
    zamiast ręcznie alokowanych tablic.
    \item Zasoby systemowe obejmuj RAII i odpowiednim deleterem;
    osobno analizuj bezpieczeństwo wątkowe wskazywanego obiektu.
\end{itemize}

\paragraph{Konkluzja.}
Najważniejsze pytanie brzmi nie ``jak zdobyć adres?'', ale
``\emph{kto jest właścicielem obiektu, jak długo on istnieje
i kto ma prawo z niego korzystać?}''. Właśnie te trzy kwestie
rozstrzygają, czy w danym miejscu potrzebujemy referencji,
surowego wskaźnika, \texttt{unique\_ptr}, \texttt{shared\_ptr}
czy \texttt{weak\_ptr}.

% Dalsza lektura:
% https://en.cppreference.com/cpp/memory/unique_ptr
% https://en.cppreference.com/cpp/memory/shared_ptr
% https://en.cppreference.com/cpp/memory/weak_ptr
% https://en.cppreference.com/cpp/memory/enable_shared_from_this
% https://en.cppreference.com/cpp/container/span


\chapter[Wskaźniki i inteligentne wskaźniki w C++]{Wskaźniki, zarządzanie pamięcią i inteligentne wskaźniki w C++}
\label{chap:wskazniki-cpp}

\section{Po co istnieją wskaźniki?}

Wskaźnik (ang. \emph{pointer}) przechowuje wartość pozwalającą
odwołać się do obiektu albo funkcji. Typowy wskaźnik do obiektu ma
postać \texttt{T*}. W odróżnieniu od obiektu typu \texttt{T},
wskaźnik może wskazywać na inny obiekt, może być pusty, a samo
wskazywanie \emph{nie oznacza posiadania wskazywanego obiektu}.

Wskaźniki wykorzystuje się między innymi do:
\begin{itemize}
    \item przekazywania i modyfikowania istniejących obiektów bez ich kopiowania;
    \item opisywania opcjonalnej relacji ``obiekt istnieje albo go nie ma'';
    \item implementowania dynamicznych struktur danych, takich jak listy,
    drzewa i grafy;
    \item realizowania polimorfizmu i dynamicznego wiązania;
    \item współpracy z interfejsami C, sterownikami i bibliotekami systemowymi;
    \item kontrolowania czasu życia zasobów za pomocą
    \texttt{std::unique\_ptr}, \texttt{std::shared\_ptr} i
    \texttt{std::weak\_ptr}.
\end{itemize}

\paragraph{Najważniejsza zasada.}
W nowoczesnym C++ wskaźnik \texttt{T*} traktujemy przede wszystkim
jako \emph{nieposiadający dostęp} do obiektu (\emph{observer}).
Odpowiedzialność za czas życia obiektu powinna być reprezentowana
osobno i jednoznacznie. Najczęściej wystarcza obiekt przechowywany
bezpośrednio, kontener standardowy albo \texttt{std::unique\_ptr}.

\section{Podstawy: adres, dereferencja i \texttt{nullptr}}

\begin{center}\small
\begin{tabularx}{\textwidth}{@{}p{0.32\textwidth}X@{}}
\toprule
Zapis & Za co odpowiada \\
\midrule
\texttt{\&x} & Pobiera adres obiektu \code{x}. \\
\code{*p} & Uzyskuje obiekt wskazywany przez \code{p}; wymaga ważnego, niepustego wskaźnika. \\
\code{p->field} & Udostępnia składową przez wskaźnik; odpowiada \code{(*p).field}. \\
\code{nullptr} & Oznacza brak wskazywanego obiektu. \\
\code{std::span<T>} & Łączy dostęp do kolejnych elementów z ich liczbą, ale nie posiada pamięci. \\
\code{std::addressof(x)} & Pobiera rzeczywisty adres obiektu nawet przy przeciążonym operatorze \texttt{\&}. \\
\bottomrule
\end{tabularx}
\end{center}

\subsection{Deklarowanie wskaźników}

Operator \texttt{\&} zastosowany do obiektu daje jego adres, natomiast
operator jednoargumentowy \texttt{*} służy do dereferencji (dostępu do
wskazywanego obiektu):

\begin{lstlisting}
#include <iostream>

int main() {
    int x = 10;
    int* p = &x;             // p points to x

    std::cout << x << '\n';  // 10
    std::cout << *p << '\n'; // 10

    *p = 25;                 // modifies x
    std::cout << x << '\n';  // 25
}
\end{lstlisting}

W deklaracji \texttt{int* p} gwiazdka należy do deklaratora, nie jest
operacją dereferencji. Nie należy też utożsamiać wskaźnika z samym
obiektem: \texttt{p} jest osobnym obiektem o własnej wartości i adresie.

\begin{lstlisting}
int x = 7;
int* p = &x;
int** pp = &p;  // pointer to a pointer

**pp = 9;       // changes x
\end{lstlisting}

Zapis \texttt{int* a, b;} deklaruje wskaźnik \texttt{a} oraz
zwykłą zmienną \texttt{b}. Dla czytelności lepiej deklarować je
oddzielnie.

\subsection{Wskaźnik pusty i poprawna dereferencja}

Pusty wskaźnik zapisujemy za pomocą \texttt{nullptr} (od C++11).
Nie wolno go dereferencjonować. Samo porównanie z
\texttt{nullptr} jest bezpieczne:

\begin{lstlisting}
int* p = nullptr;

if (p != nullptr) {
    *p = 5;
}

if (p) {                 // equivalent non-null test
    *p = 6;
}
\end{lstlisting}

\texttt{nullptr} ma osobny typ \texttt{std::nullptr\_t}; w odróżnieniu od
historycznych zapisów \texttt{0} i \texttt{NULL} nie jest zwykłą
liczbą całkowitą. Pomaga to przy przeciążaniu funkcji:

\begin{lstlisting}
void f(int);
void f(int*);

f(0);         // calls f(int)
f(nullptr);   // calls f(int*)
\end{lstlisting}

\paragraph{Uwaga o wskaźnikach niezainicjalizowanych.}
Lokalne \texttt{int* p;} bez inicjalizacji ma nieokreśloną wartość;
korzystanie z niej może prowadzić do niezdefiniowanego zachowania.
Zawsze inicjalizuj wskaźnik poprawnym adresem lub \texttt{nullptr}.

\subsection{Dostęp do składowych: \texttt{.} i \texttt{->}}

Operator \texttt{.} wybiera składową obiektu, natomiast
\texttt{->} wybiera składową przez wskaźnik:

\begin{lstlisting}
#include <string>

struct User {
    std::string name;
    int score = 0;
};

User u{"Anna", 12};
User* p = &u;

u.score = 13;
p->score = 14;       // same as (*p).score = 14
\end{lstlisting}

Nawiasy w \texttt{(*p).score} są potrzebne ze względu na priorytety
operatorów. Zapis \texttt{*p.score} nie oznacza tego samego.

\section{Wskaźniki, referencje i kwalifikator \texttt{const}}

\subsection{Kiedy wskaźnik, a kiedy referencja?}

Referencja \texttt{T\&} jest zwykle odpowiednia, gdy obiekt musi
istnieć i parametr nie jest opcjonalny. Wskaźnik \texttt{T*} nadaje
się, gdy brak obiektu jest dopuszczalny lub gdy chcemy zmieniać
obiekt, na który wskazujemy.

\begin{lstlisting}
void increase(int& x) { ++x; }

void increase_if_present(int* x) {
    if (x) { ++*x; }
}

int n = 3;
increase(n);             // n == 4
increase_if_present(&n); // n == 5
increase_if_present(nullptr); // no operation
\end{lstlisting}

Referencja również może stać się \emph{wisząca}, gdy obiekt,
do którego się odnosi, przestanie istnieć. Typ \texttt{T\&}
nie gwarantuje samodzielnie, że czas życia obiektu jest wystarczająco długi.

\subsection{Co jest stałe: wskaźnik czy obiekt?}

\begin{lstlisting}
int a = 1;
int b = 2;

const int* p = &a;       // read-only access through p
p = &b;                 // OK: p can point elsewhere
// *p = 3;              // error

int* const q = &a;       // q cannot be reassigned
*q = 3;                 // OK: a is mutable
// q = &b;              // error

const int* const r = &a; // neither reassignment nor write
int const* s = &a;       // same type as const int*
\end{lstlisting}

Reguła praktyczna: \texttt{const} po \texttt{*} dotyczy samego
wskaźnika, a przed \texttt{*} --- wskazywanego typu.

\paragraph{Pozorna niemutowalność.}
Wskaźnik \texttt{const T*} nie pozwala modyfikować obiektu
\emph{przez ten wskaźnik}. Nie oznacza, że obiekt nie może ulec zmianie
inną drogą:

\begin{lstlisting}
int x = 10;
const int* read_only = &x;
x = 20;                 // allowed
// *read_only = 30;     // not allowed
\end{lstlisting}

\section{Czas życia obiektu i wskaźniki wiszące}

\subsection{Adres nie oznacza prawa dostępu}

Wskaźnik może przechowywać adres miejsca, w którym obiekt już
nie istnieje. Jest wtedy wskaźnikiem \emph{wiszącym}
(ang. \emph{dangling pointer}). Dereferencja takiego wskaźnika
prowadzi do niezdefiniowanego zachowania (\emph{undefined behavior}, UB).

\begin{lstlisting}
int* bad() {
    int local = 42;
    return &local;       // WRONG: local dies at return
}

// int x = *bad();       // UB: dangling pointer
\end{lstlisting}

Prawidłowo można zwrócić samą wartość:

\begin{lstlisting}
int good() {
    int local = 42;
    return local;        // safe value return
}
\end{lstlisting}

\subsection{Unieważnianie wskaźników przez kontenery}

Nie każdy wskaźnik staje się wiszący przez wyjście z funkcji.
Relokacja elementów \texttt{std::vector} także może unieważnić
wskaźniki, referencje i iteratory do tych elementów:

\begin{lstlisting}
#include <vector>

std::vector<int> values{10, 20};
int* p = &values[0];

values.reserve(100);     // may reallocate storage
// *p = 50;              // unsafe if reallocation occurred
p = &values[0];          // reacquire a valid pointer
\end{lstlisting}

Po zmianie pojemności wymagającej realokacji adresy elementów
się zmieniają. W ogólności po operacji mogącej unieważnić wskaźnik
należy pobrać go ponownie, a nie próbować odgadnąć, czy nadal działa.

\subsection{Trzy niezależne pytania}

Przy każdym wskaźniku rozstrzygaj:
\begin{enumerate}
    \item Czy wartość wskaźnika jest poprawna (np. nie jest
    niezainicjalizowana)?
    \item Czy obiekt, na który wskazujemy, nadal istnieje?
    \item Czy operacja jest dozwolona dla danego typu i zakresu pamięci
    (np. czy nie wychodzi poza tablicę)?
\end{enumerate}

Sam test \texttt{p != nullptr} rozwiązuje co najwyżej fragment
pierwszego problemu, ale \emph{nie} wykrywa wiszących wskaźników.

\section{Tablice, arytmetyka wskaźników i \texttt{std::span}}

\subsection{Arytmetyka wskaźników}

Dla wskaźnika do elementu tablicy operacje \texttt{p + k} i
\texttt{p[k]} odnoszą się do jej elementów, pod warunkiem że
poruszamy się w dozwolonym zakresie. Można utworzyć wskaźnik
\emph{jeden za końcem} tablicy, ale nie wolno go dereferencjonować.

\begin{lstlisting}
int a[]{4, 8, 12};
int* p = a;          // points to a[0]

int x = *(p + 1);    // 8
int y = p[2];       // 12
int* end = a + 3;   // valid one-past-the-end pointer
// int z = *end;     // UB
\end{lstlisting}

Dodawanie i odejmowanie wskaźników ma sens w obrębie \emph{tej samej}
tablicy (z uwzględnieniem pozycji jeden za końcem). Nie należy
wykorzystywać takich działań do przemieszczania się między dowolnymi
obiektami w pamięci.

\subsection{Tablica to nie wskaźnik}

Typ \texttt{int[3]} i typ \texttt{int*} są różne, choć w wielu
wyrażeniach tablica jest automatycznie konwertowana na wskaźnik
do pierwszego elementu (tzw. \emph{array-to-pointer decay}).

\begin{lstlisting}
int a[3] = {1, 2, 3};
int* p = a;

static_assert(sizeof(a) == 3 * sizeof(int));
// sizeof(p) is the size of a pointer, not of the array
\end{lstlisting}

Funkcja \texttt{void f(int* p)} nie zna długości tablicy.
Jeśli trzeba pracować z ciągłym fragmentem pamięci o znanym
rozmiarze, C++20 udostępnia \texttt{std::span}:

\begin{lstlisting}
#include <span>
#include <vector>

int sum(std::span<const int> xs) {
    int result = 0;
    for (int x : xs) {
        result += x;
    }
    return result;
}

int a[]{1, 2, 3};
std::vector<int> v{4, 5, 6};

int s1 = sum(a);  // 6
int s2 = sum(v);  // 15
\end{lstlisting}

\texttt{std::span<T>} przechowuje widok ciągłych danych
(z informacją o liczbie elementów), ale \emph{nie jest właścicielem}
pamięci. Gdy dane znikną albo przeniosą się w inne miejsce,
\texttt{span} również może stać się niepoprawny. Dla większości
dynamicznych tablic własnościowych preferuj \texttt{std::vector<T>}.

\section{\texttt{void*}, konwersje i rzutowania wskaźników}

\subsection{\texttt{void*}: adres bez informacji o typie}

Wskaźnik \texttt{void*} może przechowywać adres obiektu bez
informacji o jego konkretnym typie. Bywa używany w interfejsach C
i bibliotekach niskopoziomowych, ale nie pozwala na bezpośrednią
dereferencję:

\begin{lstlisting}
int x = 42;
void* opaque = &x;              // type erased
// *opaque = 7;                 // error: cannot dereference void*

int* p = static_cast<int*>(opaque);
*p = 7;                        // valid: original object is int
\end{lstlisting}

Przywrócenie błędnego typu i późniejsza próba użycia obiektu
mogą prowadzić do niezdefiniowanego zachowania. Jeśli typ jest
znany, zwykle lepiej przechowywać od razu \texttt{T*}.

\subsection{\texttt{static\_cast}, \texttt{dynamic\_cast} i
\texttt{reinterpret\_cast}}

\begin{itemize}
    \item \texttt{static\_cast<T*>} nadaje się m.in. do poprawnych
    konwersji pomiędzy \texttt{void*} a wskaźnikiem obiektowym,
    jeśli rzeczywisty typ i wyrównanie na to pozwalają.
    \item \texttt{dynamic\_cast<T*>} służy do sprawdzanych w czasie
    wykonania konwersji w hierarchiach klas polimorficznych.
    Niepowodzenie konwersji wskaźnika zwraca \texttt{nullptr}.
    \item \texttt{reinterpret\_cast<T*>} umożliwia niskopoziomową
    reinterpretację reprezentacji wskaźnika, ale sama konwersja
    nie uprawnia do dowolnej dereferencji ani nie tworzy obiektu
    odpowiedniego typu.
\end{itemize}

\begin{lstlisting}
struct Base {
    virtual ~Base() = default;
};

struct Derived : Base {
    int value = 12;
};

Derived d;
Base* base = &d;

if (Derived* p = dynamic_cast<Derived*>(base)) {
    p->value = 30;          // checked downcast succeeded
}
\end{lstlisting}

Przy przekazywaniu parametrów funkcji nie używaj rzutowań
jako sposobu na obejście \texttt{const}, sprawdzania czasu życia
ani zasad własności. Dobrze zaprojektowany interfejs zazwyczaj
w ogóle nie wymaga rzutowania wskaźników.

\section{Pamięć dynamiczna: \texttt{new}, \texttt{delete} i ich ograniczenia}

\subsection{Klasyczny mechanizm}

Operator \texttt{new} tworzy obiekt w pamięci przydzielonej
dynamicznie i zwraca wskaźnik. Odpowiadający mu
\texttt{delete} niszczy obiekt i zwalnia pamięć:

\begin{lstlisting}
int* p = new int(42);
// use *p

delete p;
p = nullptr;           // avoids this one stale alias
\end{lstlisting}

Dla tablic obowiązuje para \texttt{new[]} i \texttt{delete[]}:

\begin{lstlisting}
int* a = new int[5]{}; // zero-initialized elements
// use a[0], ..., a[4]

delete[] a;
\end{lstlisting}

\textbf{Nie wolno} łączyć \texttt{new[]} z \texttt{delete} ani
zwykłego \texttt{new} z \texttt{delete[]}. Ponowne
\texttt{delete} tego samego żywego przydziału również jest błędem.
\texttt{delete nullptr;} jest natomiast bezpieczne.

\subsection{Dlaczego ręczne zarządzanie pamięcią jest niebezpieczne?}

\begin{lstlisting}
void fragile() {
    int* p = new int(42);
    operation_that_may_throw();
    delete p;   // skipped if the operation throws
}
\end{lstlisting}

Jeśli funkcja zgłosi wyjątek, pamięć zostaje utracona
(\emph{memory leak}). Analogiczny problem powstaje przy
wielu instrukcjach \texttt{return}, kopiowaniu wskaźników i
niespójnym przekazywaniu odpowiedzialności za \texttt{delete}.

\paragraph{Zalecenie.}
W typowym kodzie aplikacyjnym nie pisz bezpośrednio
\texttt{new}/\texttt{delete}. Użyj obiektu lokalnego,
\texttt{std::vector<T>} albo inteligentnego wskaźnika.
Ręczne operacje alokowania pozostają potrzebne w implementacjach
niskopoziomowych, alokatorach i przy integracji z określonym API.

\section{RAII i semantyka własności}

\subsection{RAII: zasób należy do obiektu}

\emph{Resource Acquisition Is Initialization} (RAII) oznacza
powiązanie przejęcia zasobu z utworzeniem obiektu zarządzającego,
a zwolnienia zasobu z jego destruktorem. Stosuje się je nie tylko
do pamięci, ale też do plików, blokad, gniazd sieciowych i uchwytów
systemowych.

\begin{lstlisting}
#include <fstream>
#include <string>

void save(const std::string& path) {
    std::ofstream file(path);
    file << "example\n";
} // file is closed by its destructor
\end{lstlisting}

Destruktor uruchamia się podczas normalnego opuszczania zakresu
oraz podczas odwijania stosu po wyjątku. Stąd RAII upraszcza
bezpieczeństwo wyjątkowe.

\subsection{Trzy rodzaje relacji z obiektem}

\begin{description}
    \item[Własność wyłączna (\emph{unique ownership}).]
    Dokładnie jeden obiekt zarządzający odpowiada za zniszczenie
    zasobu: najczęściej \texttt{std::unique\_ptr<T>}.
    \item[Własność współdzielona (\emph{shared ownership}).]
    Wielu właścicieli może utrzymywać jeden zasób przy życiu:
    \texttt{std::shared\_ptr<T>}.
    \item[Dostęp bez własności (\emph{non-owning access}).]
    Użytkownik widzi obiekt, ale nie przedłuża jego życia:
    zwykły wskaźnik, referencja, \texttt{std::span<T>} lub
    \texttt{std::weak\_ptr<T>} (ten ostatni obserwuje stan
    zarządzany przez \texttt{shared\_ptr} i pozwala warunkowo
    uzyskać właściciela).
\end{description}

Dobór reprezentacji wynika z \emph{kontraktu własności}, a nie
z samego faktu, że obiekt znajduje się w pamięci dynamicznej.

\section{\texttt{std::unique\_ptr}: jeden właściciel}

\begin{center}\small
\begin{tabularx}{\textwidth}{@{}p{0.36\textwidth}X@{}}
\toprule
Mechanizm & Własność i czas życia \\
\midrule
\code{std::make_unique<T>(...)} & Tworzy obiekt i zwraca jego jedynego właściciela. \\
\code{unique_ptr::get()} & Zwraca surowy wskaźnik bez przekazania własności. \\
\code{unique_ptr::reset(p)} & Niszczy obecnie posiadany obiekt i przejmuje wskaźnik \code{p}. \\
\code{unique_ptr::release()} & Oddaje surowy wskaźnik bez zniszczenia obiektu; odtąd trzeba osobno zadbać o jego zwolnienie. \\
\code{std::shared_ptr<T>} & Współdzieli własność; obiekt znika po odejściu ostatniego właściciela. \\
\code{weak_ptr::lock()} & Próbuje uzyskać \code{shared_ptr}; zwraca pusty wynik, jeśli obiekt wygasł. \\
\bottomrule
\end{tabularx}
\end{center}

\subsection{Tworzenie obiektu i automatyczne niszczenie}

\texttt{std::unique\_ptr<T>} (nagłówek \texttt{<memory>}, C++11)
zarządza pojedynczym obiektem przez wyłączną własność. Od C++14
najczęściej tworzymy go przez \texttt{std::make\_unique<T>}:

\begin{lstlisting}
#include <memory>
#include <string>

#include <utility>

struct Person {
    explicit Person(std::string n) : name(std::move(n)) {}
    std::string name;
};

void example() {
    auto p = std::make_unique<Person>("Ala");
    p->name = "Ola";
    // no delete needed
} // destructor destroys Person
\end{lstlisting}

\texttt{unique\_ptr} można traktować jako obiekt zarządzający:
jego destruktor wywołuje odpowiednią funkcję usuwającą zasób
(\emph{deleter}). Domyślnie jest nią \texttt{std::default\_delete<T>}.

\subsection{Brak kopiowania i przenoszenie własności}

Wskaźnik o własności wyłącznej nie może być kopiowany.
Może natomiast zostać \emph{przeniesiony} za pomocą
\texttt{std::move}:

\begin{lstlisting}
#include <memory>
#include <utility>

auto first = std::make_unique<int>(10);
// auto bad = first;                  // error: copy

auto second = std::move(first);      // transfers ownership
// first == nullptr
*second = 20;
\end{lstlisting}

\texttt{std::move} samo niczego nie zwalnia i nie przenosi fizycznie
wskazywanego obiektu. Zwraca wyrażenie pozwalające wywołać konstruktor
lub przypisanie przenoszące. W tym przypadku przenoszona jest
\emph{własność wskaźnika}, a obiekt typu \texttt{int} nie zmienia adresu.

\subsection{Przekazywanie do funkcji: kto jest właścicielem?}

Rozróżnijmy trzy interfejsy:

\begin{lstlisting}
#include <memory>
#include <utility>

struct Widget { void run() {} };

void inspect(const Widget& w);             // borrows, non-null
void maybe_inspect(const Widget* w);       // borrows, nullable
void take(std::unique_ptr<Widget> w);      // takes ownership

void demo() {
    auto p = std::make_unique<Widget>();
    inspect(*p);                // p still owns the object
    maybe_inspect(p.get());     // p still owns the object
    take(std::move(p));         // ownership transferred
    // p is now empty
}
\end{lstlisting}

Jeżeli funkcja \emph{jedynie używa} obiektu, przekazywanie przez
\texttt{T\&}, \texttt{const T\&} lub \texttt{T*} jest zwykle
czytelniejsze niż przekazywanie \texttt{unique\_ptr} przez referencję.
Parametr \texttt{unique\_ptr<T>} przyjmowany \emph{przez wartość}
sygnalizuje przejęcie własności. Dla funkcji, która ma zastąpić
obiekt będący własnością wywołującego, można użyć
\texttt{std::unique\_ptr<T>\&}.

\subsection{Metody \texttt{get}, \texttt{reset}, \texttt{release}, \texttt{swap}}

\begin{lstlisting}
#include <memory>
#include <utility>

auto p = std::make_unique<int>(7);
int* observer = p.get();       // borrow; do not delete

p.reset(new int(9));           // destroys old int
// observer is now dangling

auto q = std::make_unique<int>(11);
p.swap(q);                    // exchanges managed objects

int* raw = q.release();        // q becomes empty; raw owns now
// responsibility for delete raw is now explicit
delete raw;
\end{lstlisting}

\texttt{get()} zwraca dostęp bez własności. \texttt{reset()} zastępuje
zarządzany obiekt, zwalniając dotychczasowy. \texttt{release()} oddaje
surowy wskaźnik \emph{bez zwolnienia obiektu}; dlatego jest potrzebne
głównie przy przekazywaniu własności do interfejsu, który tego wymaga.
Nie należy mylić \texttt{release()} z \texttt{reset()}.

\subsection{Tablice: \texttt{unique\_ptr<T[]>} a \texttt{vector<T>}}

\begin{lstlisting}
#include <memory>
#include <vector>

auto a = std::make_unique<int[]>(5); // array of 5 zeros
a[2] = 17;

std::vector<int> b(5);         // usually more convenient
b.at(2) = 17;                  // checked indexing
\end{lstlisting}

\texttt{unique\_ptr<T[]>} dobiera \texttt{delete[]} i udostępnia
\texttt{operator[]}, lecz sam nie przechowuje rozmiaru.
Zwykle preferuj \texttt{std::vector<T>}, bo posiada rozmiar,
iteratory i bogatszy interfejs. Dla tablicy o stałym rozmiarze
określonym w kodzie rozważ \texttt{std::array<T, N>}.

\subsection{Własny deleter: zasób inny niż pamięć \texttt{new}}

Inteligentny wskaźnik może wywoływać dowolną funkcję zwalniającą
zasób. Typowym przykładem jest plik w interfejsie C:

\begin{lstlisting}
#include <cstdio>
#include <memory>

struct FileCloser {
    void operator()(std::FILE* file) const noexcept {
        if (file) { std::fclose(file); }
    }
};

using FilePtr = std::unique_ptr<std::FILE, FileCloser>;

void example() {
    FilePtr file(std::fopen("log.txt", "w"));
    if (!file) { return; }

    std::fputs("hello\n", file.get());
} // fclose is called automatically
\end{lstlisting}

W praktyce deleter bywa osobną strukturą, a jego typ wchodzi
w skład typu \texttt{unique\_ptr<T, Deleter>}. Mechanizm pozwala
owinąć RAII również zasoby zwalniane przez \texttt{free()},
\texttt{fclose()} lub bibliotekę zewnętrzną.

\subsection{Polimorfizm i wirtualny destruktor}

\begin{lstlisting}
#include <memory>

struct Shape {
    virtual ~Shape() = default;
    virtual double area() const = 0;
};

struct Circle : Shape {
    explicit Circle(double r) : radius(r) {}
    double area() const override { return 3.14159 * radius * radius; }
    double radius;
};

std::unique_ptr<Shape> make_shape() {
    return std::make_unique<Circle>(2.0);
}
\end{lstlisting}

Przy niszczeniu obiektu pochodnego przez wskaźnik do bazy
zwykle potrzebny jest wirtualny destruktor bazy.
Dotyczy to również przenoszenia własności z klasy pochodnej
do bazowej przy użyciu \texttt{unique\_ptr} i domyślnego deletera.

\section{\texttt{std::shared\_ptr}: współdzielona własność}

\subsection{Jak działa licznik właścicieli?}

\texttt{std::shared\_ptr<T>} (C++11) udostępnia współdzieloną
własność zasobu. Zwykle tworzymy go za pomocą
\texttt{std::make\_shared<T>}:

\begin{lstlisting}
#include <memory>

auto a = std::make_shared<int>(10);
auto b = a;             // shared ownership

// a.use_count() == 2 (provided there are no other owners)
b.reset();              // releases b's ownership
// a.use_count() == 1
\end{lstlisting}

Gdy liczba \emph{silnych właścicieli} spadnie do zera, zarządzany
obiekt zostaje zniszczony. Implementacja utrzymuje tzw.
\emph{blok kontrolny} (\emph{control block}) zawierający między
innymi liczniki właścicieli, informacje potrzebne do zniszczenia
zasobu i ewentualny deleter.

\paragraph{Schemat ideowy.}
\begin{lstlisting}
shared_ptr a ---\
                +--> [ control block: strong = 2 ] --> [ T ]
shared_ptr b ---/           weak = 0
\end{lstlisting}

\texttt{use\_count()} ma znaczenie przede wszystkim diagnostyczne.
Nie powinno się na podstawie odczytanego licznika podejmować decyzji
o bezpieczeństwie wątkowym ani o wyłącznej własności, która może
zmienić się równolegle.

\subsection{Tworzenie i przekazywanie \texttt{shared\_ptr}}

\begin{lstlisting}
#include <memory>

struct Data { int x = 0; };

std::shared_ptr<Data> create() {
    return std::make_shared<Data>();
}

void keep(std::shared_ptr<Data> p) {
    // stores a copy of p somewhere, if needed
}

void read_only(const Data& data) {
    // does not require shared ownership
}

void demo() {
    auto p = create();
    read_only(*p); // borrow
    keep(p);       // creates another temporary shared owner
}
\end{lstlisting}

Parametr \texttt{shared\_ptr<T>} przyjmowany przez wartość oznacza,
że funkcja otrzymuje własny udział we własności na czas wywołania
i ewentualnie może go przechować. Jeśli funkcja wyłącznie
czyta obiekt, zwykle wystarcza \texttt{const T\&} lub \texttt{T*}.
Samo przekazanie \texttt{const shared\_ptr<T>\&} unika inkrementacji
licznika, ale bardziej wiąże interfejs z modelem własności.

\subsection{\texttt{make\_shared} a konstruktor ze wskaźnika}

W typowych implementacjach \texttt{make\_shared} umieszcza obiekt
i blok kontrolny w jednym przydziale pamięci, co może ograniczać
narzut alokacji. Natomiast konstrukcja \texttt{shared\_ptr<T>(new T)}
zwykle wymaga osobnych przydziałów.

\begin{lstlisting}
auto preferred = std::make_shared<int>(42);
std::shared_ptr<int> also_valid(new int(42));
\end{lstlisting}

Nie wolno jednak tworzyć dwóch \texttt{shared\_ptr} niezależnie
z \emph{tego samego} surowego wskaźnika:

\begin{lstlisting}
int* raw = new int(42);
std::shared_ptr<int> a(raw);
// std::shared_ptr<int> b(raw); // WRONG: second control block
\end{lstlisting}

Dwa niezależne bloki kontrolne próbowałyby zwolnić ten sam obiekt.
Poprawne współdzielenie powstaje przez \emph{kopiowanie istniejącego}
\texttt{shared\_ptr}: \texttt{auto b = a;}.

\paragraph{Kiedy nie stosować \texttt{make\_shared}?}
Zintegrowany przydział może zatrzymywać pamięć zajmowaną przez
obiekt w przydzielonym bloku do momentu zniszczenia także ostatniego
\texttt{weak\_ptr}, nawet gdy destruktor obiektu już się wykonał.
Przy dużych obiektach i długo żyjących słabych referencjach
warto rozważyć osobny przydział. Innymi powodami mogą być
specjalne wymagania alokacji lub konstruktory niedostępne
dla \texttt{make\_shared}.

\subsection{Współdzielenie własności a współdzielenie mutacji}

Sam \texttt{shared\_ptr} nie synchronizuje dostępu do pól obiektu:

\begin{lstlisting}
#include <memory>
#include <mutex>

struct Counter {
    std::mutex mutex;
    int value = 0;

    void increment() {
        std::lock_guard<std::mutex> lock(mutex);
        ++value;
    }
};

auto p = std::make_shared<Counter>();
// many threads may own copies of p, but must synchronize
// concurrent modifications of Counter::value
\end{lstlisting}

Blok kontrolny pozwala różnym egzemplarzom \texttt{shared\_ptr}
bezpiecznie współdzielić własność. Równoczesna modyfikacja
\emph{tego samego} egzemplarza wymaga jednak synchronizacji.
Od C++20 można użyć \texttt{std::atomic<std::shared\_ptr<T>>}.

\section{\texttt{std::weak\_ptr}: obserwacja bez utrzymywania przy życiu}

\subsection{Problem cykli własności}

Gdy obiekty wzajemnie posiadają się przez \texttt{shared\_ptr},
liczniki właścicieli mogą nigdy nie spaść do zera:

\begin{lstlisting}
#include <memory>

struct Node {
    std::shared_ptr<Node> next;
};

void leak_cycle() {
    auto a = std::make_shared<Node>();
    auto b = std::make_shared<Node>();
    a->next = b;
    b->next = a;          // cycle of strong ownership
} // nodes remain alive after both locals are destroyed
\end{lstlisting}

To nie jest problem \emph{niezerowego wskaźnika}, lecz błędnego
modelu własności: każdy węzeł utrzymuje przy życiu drugi.

\subsection{Przerywanie cyklu i \texttt{lock()}}

\texttt{std::weak\_ptr<T>} obserwuje ten sam blok kontrolny,
ale nie zwiększa liczby silnych właścicieli:

\begin{lstlisting}
#include <memory>

struct Node {
    std::weak_ptr<Node> parent;         // non-owning
    std::shared_ptr<Node> child;        // owning
};

void example() {
    auto root = std::make_shared<Node>();
    auto leaf = std::make_shared<Node>();

    root->child = leaf;
    leaf->parent = root;

    if (auto parent = leaf->parent.lock()) {
        // parent is a temporary shared owner
        (void)parent;
    }
} // no ownership cycle
\end{lstlisting}

Metoda \texttt{lock()} atomowo próbuje uzyskać
\texttt{shared\_ptr<T>}. Jeżeli obiekt jeszcze istnieje,
zwraca właściciela utrzymującego go przy życiu przez czas
używania. W przeciwnym razie zwraca pusty \texttt{shared\_ptr}.

\begin{lstlisting}
std::weak_ptr<int> observer;
{
    auto p = std::make_shared<int>(8);
    observer = p;
    if (auto alive = observer.lock()) {
        *alive += 1;           // safe while alive owns it
    }
}

// observer.expired() is now true
// observer.lock() now returns an empty shared_ptr
\end{lstlisting}

W kodzie współbieżnym nie stosuj schematu
\texttt{if (!w.expired()) { ... }} z późniejszym odwołaniem:
między sprawdzeniem a użyciem obiekt może zniknąć.
Preferuj jedną operację \texttt{lock()} i posługuj się
otrzymanym \texttt{shared\_ptr}.

\section{\texttt{enable\_shared\_from\_this} i wskaźnik \texttt{this}}

Gdy metoda obiektu zarządzanego przez \texttt{shared\_ptr}
chce przekazać dalej \emph{współdzieloną własność do siebie},
nie może wykonać \texttt{std::shared\_ptr<T>(this)}:
powstałby drugi, niezależny blok kontrolny.

Można zamiast tego dziedziczyć publicznie po
\texttt{std::enable\_shared\_from\_this<T>}:

\begin{lstlisting}
#include <memory>

struct Session : std::enable_shared_from_this<Session> {
    std::shared_ptr<Session> self() {
        return shared_from_this();
    }
};

void example() {
    auto session = std::make_shared<Session>();
    auto second_owner = session->self();
    // both share the same control block
}
\end{lstlisting}

Metodę \texttt{shared\_from\_this()} wolno wywołać dopiero wtedy,
gdy obiektem zarządza już \texttt{shared\_ptr}. Jej wywołanie
na obiekcie lokalnym albo w konstruktorze, zanim powstanie
właściciel, zgłasza wyjątek \texttt{std::bad\_weak\_ptr}.

Od C++17 dostępna jest również metoda \texttt{weak\_from\_this()}.

\section{Zaawansowany przykład: wskaźnik aliasujący \texttt{shared\_ptr}}

\texttt{shared\_ptr} rozróżnia \emph{obiekt zarządzany} (którego
czas życia kontroluje blok kontrolny) i \emph{przechowywany
wskaźnik} (zwracany przez \texttt{get()}). Specjalny konstruktor
aliasujący pozwala udostępnić adres składowej, jednocześnie
utrzymując przy życiu cały obiekt:

\begin{lstlisting}
#include <memory>

struct Pair {
    int first = 1;
    int second = 2;
};

void example() {
    auto whole = std::make_shared<Pair>();
    std::shared_ptr<int> member(whole, &whole->second);

    whole.reset();       // Pair is still alive
    *member = 20;        // accesses Pair::second safely
} // member destroys the last strong owner
\end{lstlisting}

Konstruktor aliasujący jest przydatny np. przy udostępnianiu
podobiektów, ale wymaga prawidłowej relacji czasów życia.
Nie należy go traktować jako sposobu na naprawienie
wcześniej wiszącego wskaźnika.

\section{Wskaźniki w interfejsach funkcji i klas}

\subsection{Katalog typowych kontraktów}

\begin{center}
\small
\begin{tabular}{|p{0.31\textwidth}|p{0.56\textwidth}|}
\hline
\textbf{Interfejs} & \textbf{Znaczenie} \\
\hline
\texttt{void f(const T\&)} & Obiekt wymagany, tylko odczyt; brak przekazania własności. \\
\hline
\texttt{void f(T\&)} & Obiekt wymagany, możliwa modyfikacja; brak przekazania własności. \\
\hline
\texttt{void f(const T*)} & Obiekt opcjonalny, tylko odczyt; brak własności. \\
\hline
\texttt{void f(T*)} & Obiekt opcjonalny, możliwa modyfikacja; brak własności. \\
\hline
\texttt{void f(unique\_ptr<T>)} & Przejęcie wyłącznej własności. \\
\hline
\texttt{unique\_ptr<T> f()} & Przekazanie wyłącznej własności wywołującemu. \\
\hline
\texttt{void f(shared\_ptr<T>)} & Uzyskanie udziału we współdzielonej własności. \\
\hline
\texttt{void f(weak\_ptr<T>)} & Możliwość obserwacji obiektu, którego życie kontrolują inni. \\
\hline
\texttt{void f(span<const T>)} & Widok na ciągłą sekwencję, bez własności. \\
\hline
\end{tabular}
\end{center}

W tabeli pominięto przestrzeń nazw \texttt{std::} dla czytelności.
To \emph{konwencje projektowania API}, a nie automatyczne
gwarancje długowieczności surowych wskaźników i referencji.

\subsection{Dlaczego nie zawsze używać \texttt{shared\_ptr}?}

Rozważmy funkcję, która wyświetla nazwę:

\begin{lstlisting}
#include <iostream>
#include <memory>
#include <string>

struct Product { std::string name; };

void print(const Product& p) {
    std::cout << p.name << '\n';
}

void example() {
    Product local{"book"};
    print(local);                  // no allocation

    auto owned = std::make_unique<Product>(Product{"pen"});
    print(*owned);                 // still no extra owner
}
\end{lstlisting}

Funkcja \texttt{print} nie potrzebuje znać strategii
przechowywania obiektu. Przekazanie referencji oddziela
\emph{sposób dostępu} od \emph{sposobu zarządzania czasem życia}.

\section{PImpl: ukrywanie implementacji przez \texttt{unique\_ptr}}

Technika \emph{pointer to implementation} pozwala utrzymać
szczegóły implementacji poza publicznym nagłówkiem.
Typ implementacyjny może być niekompletny w nagłówku, ale
miejsce wykonania domyślnego deletera musi widzieć jego definicję.
Dlatego destruktor klasy definiujemy w pliku \texttt{.cpp}.

\begin{lstlisting}
// widget.hpp
#pragma once
#include <memory>

class Widget {
public:
    Widget();
    ~Widget();                 // defined in .cpp

    Widget(Widget&&) noexcept;
    Widget& operator=(Widget&&) noexcept;

    Widget(const Widget&) = delete;
    Widget& operator=(const Widget&) = delete;

    void work();

private:
    struct Impl;
    std::unique_ptr<Impl> impl_;
};
\end{lstlisting}

\begin{lstlisting}
// widget.cpp
#include "widget.hpp"
#include <memory>

struct Widget::Impl {
    int value = 0;
    void run() { ++value; }
};

Widget::Widget() : impl_(std::make_unique<Impl>()) {}
Widget::~Widget() = default;
Widget::Widget(Widget&&) noexcept = default;
Widget& Widget::operator=(Widget&&) noexcept = default;

void Widget::work() { impl_->run(); }
\end{lstlisting}

Ten przykład prezentuje semantykę przenoszenia. Po przeniesieniu
źródłowy obiekt \texttt{Widget} ma pusty \texttt{impl\_};
nie wolno wywoływać na nim \texttt{work()} bez odpowiedniego
sprawdzenia. W produkcyjnym API można ten stan jawnie obsłużyć
lub określić w kontrakcie.

\section{Wskaźniki do funkcji i składowych klas}

Wskaźniki nie muszą odnosić się do danych. Można wskazywać
funkcje, a także składowe klas, choć te ostatnie wymagają
innej składni.

\subsection{Wskaźnik do funkcji}

\begin{lstlisting}
int add(int a, int b) { return a + b; }
int mul(int a, int b) { return a * b; }

using Operation = int(*)(int, int);

int apply(int x, int y, Operation op) {
    return op(x, y);
}

int a = apply(2, 3, &add); // 5
int b = apply(2, 3, &mul); // 6
\end{lstlisting}

Obiekty wywoływalne różnych typów można przechowywać za pomocą
\texttt{std::function}. Prosty szablon lub lambda bywają jednak
bardziej efektywne.

\subsection{Wskaźnik do składowej}

\begin{lstlisting}
struct Record {
    int value = 0;
    int twice() const { return 2 * value; }
};

int Record::* field = &Record::value;
int (Record::* method)() const = &Record::twice;

Record r{7};
int x = r.*field;          // 7
int y = (r.*method)();     // 14

Record* p = &r;
int z = p->*field;         // 7
\end{lstlisting}

Wskaźnik do składowej nie jest zwykłym adresem gotowego obiektu.
Dopiero w połączeniu z konkretną instancją klasy pozwala
odwołać się do odpowiedniej składowej.

\section{Wskaźniki i wielowątkowość}

\subsection{\texttt{shared\_ptr} nie eliminuje wyścigów danych}

Możliwość bezpiecznego utrzymania czasu życia obiektu nie oznacza,
że jego pola można dowolnie zmieniać z wielu wątków.
Jeśli dwa wątki zapisują do \texttt{*p}, potrzebna jest osobna
synchronizacja (np. \texttt{std::mutex} albo atomiki).

Od C++20 dostępna jest specjalizacja
\texttt{std::atomic<std::shared\_ptr<T>>}, przydatna do atomowego
publikowania nowej wersji współdzielonego obiektu:

\begin{lstlisting}
#include <atomic>
#include <memory>

struct Config {
    int workers = 4;
};

std::atomic<std::shared_ptr<const Config>> current{
    std::make_shared<const Config>(Config{4})
};

void publish(int count) {
    current.store(std::make_shared<const Config>(Config{count}));
}

int read_workers() {
    auto snapshot = current.load(); // keeps snapshot alive
    return snapshot->workers;
}
\end{lstlisting}

Tu czytelnicy otrzymują stabilny, niemutowalny \emph{snapshot}.
Atomowość dotyczy odczytu i zapisu \texttt{shared\_ptr},
a nie arbitralnych operacji na jego zawartości.

\section{Najczęstsze błędy zebrane w jednym miejscu}

\begin{enumerate}
    \item \textbf{Dereferencja \texttt{nullptr} lub wskaźnika
    niezainicjalizowanego.} Zawsze sprawdzaj kontrakt i inicjalizuj
    wskaźniki.
    \item \textbf{Wskaźnik wiszący.} \texttt{get()}, \texttt{span}
    i surowa referencja nie przedłużają życia zasobu.
    \item \textbf{Podwójne \texttt{delete}.} Jeden przydział musi mieć
    jedną, jednoznaczną strategię zwalniania.
    \item \textbf{Mieszanie \texttt{new[]} i \texttt{delete}.}
    Do tablic służy \texttt{delete[]}; zwykle lepiej wybrać kontener.
    \item \textbf{Zwracanie adresu lokalnej zmiennej.}
    Jej czas życia kończy się przy wyjściu z zakresu.
    \item \textbf{Przekazywanie surowego wskaźnika z \texttt{get()}
    do nowego \texttt{shared\_ptr}.} Tworzy niezależny blok kontrolny.
    \item \textbf{Nieuzasadnione \texttt{shared\_ptr} wszędzie.}
    Współdzielenie własności kosztuje i zaciera odpowiedzialność.
    \item \textbf{Cykl silnych właścicieli.} Przerywaj go przez
    \texttt{weak\_ptr} lub przeprojektowanie relacji.
    \item \textbf{Nadużywanie \texttt{release()}.} Po wywołaniu
    ktoś musi jawnie przejąć odpowiedzialność za zasób.
    \item \textbf{Brak wirtualnego destruktora w bazie.}
    Usuwanie obiektu pochodnego przez wskaźnik bazowy z domyślnym
    deleterem może wówczas prowadzić do UB.
    \item \textbf{Założenie, że \texttt{shared\_ptr} chroni obiekt
    przed wyścigami.} Chroni zarządzanie własnością, nie samą mutację.
    \item \textbf{Przechowywanie wskaźników do elementów \texttt{vector}
    po realokacji.} Przy operacjach zmieniających bufor pobieraj
    adresy ponownie.
\end{enumerate}

\section{Jak wybrać właściwy typ?}

\begin{center}
\small
\begin{tabular}{|p{0.43\textwidth}|p{0.43\textwidth}|}
\hline
\textbf{Sytuacja} & \textbf{Najczęstszy wybór} \\
\hline
Pojedynczy obiekt o prostym czasie życia & Zwykła zmienna \texttt{T} \\
\hline
Dynamiczna lista elementów & \texttt{std::vector<T>} \\
\hline
Obiekt wymagany tylko do odczytu & \texttt{const T\&} \\
\hline
Obiekt wymagany i modyfikowany & \texttt{T\&} \\
\hline
Obiekt opcjonalny bez przekazania własności & \texttt{T*} lub \texttt{const T*} \\
\hline
Pojedynczy właściciel dynamicznego obiektu & \texttt{std::unique\_ptr<T>} \\
\hline
Kilku niezależnych właścicieli & \texttt{std::shared\_ptr<T>} \\
\hline
Obserwacja obiektu współdzielonego & \texttt{std::weak\_ptr<T>} \\
\hline
Fragment istniejącej tablicy, bez własności & \texttt{std::span<T>} \\
\hline
Polimorficzna kolekcja obiektów & \texttt{std::vector} zawierający \texttt{unique\_ptr} \\
\hline
\end{tabular}
\end{center}

\subsection{Praktyczny przykład: kolekcja obiektów polimorficznych}

\begin{lstlisting}
#include <iostream>
#include <memory>
#include <vector>

struct Task {
    virtual ~Task() = default;
    virtual void run() const = 0;
};

struct PrintTask : Task {
    void run() const override { std::cout << "hello\n"; }
};

int main() {
    std::vector<std::unique_ptr<Task>> tasks;
    tasks.push_back(std::make_unique<PrintTask>());

    for (const auto& task : tasks) {
        task->run();
    }
} // vector and tasks destroy everything automatically
\end{lstlisting}

\section{Podsumowanie i praktyczne reguły}

\begin{itemize}
    \item Preferuj \emph{przechowywanie przez wartość}, jeśli
    obiekt nie wymaga odrębnej, dynamicznej tożsamości.
    \item Gdy obiekt ma jednego właściciela, używaj
    \texttt{std::unique\_ptr}; własność przenoś przez
    \texttt{std::move}.
    \item Sięgaj po \texttt{std::shared\_ptr} dopiero wtedy,
    gdy kilka części programu \emph{rzeczywiście} musi niezależnie
    utrzymywać ten sam obiekt przy życiu.
    \item Jeśli obserwator nie powinien utrzymywać obiektu przy życiu,
    użyj \texttt{std::weak\_ptr} i metody \texttt{lock()}.
    \item Do niewłasnościowego dostępu używaj
    \texttt{T\&}, \texttt{const T\&}, \texttt{T*} lub
    \texttt{std::span<T>} w zależności od kontraktu.
    \item Nie dereferencjonuj wskaźnika tylko dlatego, że jest
    różny od \texttt{nullptr}: obiekt może już nie istnieć.
    \item Preferuj \texttt{std::vector} i \texttt{std::array}
    zamiast ręcznie alokowanych tablic.
    \item Zasoby systemowe obejmuj RAII i odpowiednim deleterem;
    osobno analizuj bezpieczeństwo wątkowe wskazywanego obiektu.
\end{itemize}

\paragraph{Konkluzja.}
Najważniejsze pytanie brzmi nie ``jak zdobyć adres?'', ale
``\emph{kto jest właścicielem obiektu, jak długo on istnieje
i kto ma prawo z niego korzystać?}''. Właśnie te trzy kwestie
rozstrzygają, czy w danym miejscu potrzebujemy referencji,
surowego wskaźnika, \texttt{unique\_ptr}, \texttt{shared\_ptr}
czy \texttt{weak\_ptr}.

% Dalsza lektura:
% https://en.cppreference.com/cpp/memory/unique_ptr
% https://en.cppreference.com/cpp/memory/shared_ptr
% https://en.cppreference.com/cpp/memory/weak_ptr
% https://en.cppreference.com/cpp/memory/enable_shared_from_this
% https://en.cppreference.com/cpp/container/span


\chapter[Wskaźniki i inteligentne wskaźniki w C++]{Wskaźniki, zarządzanie pamięcią i inteligentne wskaźniki w C++}
\label{chap:wskazniki-cpp}

\section{Po co istnieją wskaźniki?}

Wskaźnik (ang. \emph{pointer}) przechowuje wartość pozwalającą
odwołać się do obiektu albo funkcji. Typowy wskaźnik do obiektu ma
postać \texttt{T*}. W odróżnieniu od obiektu typu \texttt{T},
wskaźnik może wskazywać na inny obiekt, może być pusty, a samo
wskazywanie \emph{nie oznacza posiadania wskazywanego obiektu}.

Wskaźniki wykorzystuje się między innymi do:
\begin{itemize}
    \item przekazywania i modyfikowania istniejących obiektów bez ich kopiowania;
    \item opisywania opcjonalnej relacji ``obiekt istnieje albo go nie ma'';
    \item implementowania dynamicznych struktur danych, takich jak listy,
    drzewa i grafy;
    \item realizowania polimorfizmu i dynamicznego wiązania;
    \item współpracy z interfejsami C, sterownikami i bibliotekami systemowymi;
    \item kontrolowania czasu życia zasobów za pomocą
    \texttt{std::unique\_ptr}, \texttt{std::shared\_ptr} i
    \texttt{std::weak\_ptr}.
\end{itemize}

\paragraph{Najważniejsza zasada.}
W nowoczesnym C++ wskaźnik \texttt{T*} traktujemy przede wszystkim
jako \emph{nieposiadający dostęp} do obiektu (\emph{observer}).
Odpowiedzialność za czas życia obiektu powinna być reprezentowana
osobno i jednoznacznie. Najczęściej wystarcza obiekt przechowywany
bezpośrednio, kontener standardowy albo \texttt{std::unique\_ptr}.

\section{Podstawy: adres, dereferencja i \texttt{nullptr}}

\begin{center}\small
\begin{tabularx}{\textwidth}{@{}p{0.32\textwidth}X@{}}
\toprule
Zapis & Za co odpowiada \\
\midrule
\texttt{\&x} & Pobiera adres obiektu \code{x}. \\
\code{*p} & Uzyskuje obiekt wskazywany przez \code{p}; wymaga ważnego, niepustego wskaźnika. \\
\code{p->field} & Udostępnia składową przez wskaźnik; odpowiada \code{(*p).field}. \\
\code{nullptr} & Oznacza brak wskazywanego obiektu. \\
\code{std::span<T>} & Łączy dostęp do kolejnych elementów z ich liczbą, ale nie posiada pamięci. \\
\code{std::addressof(x)} & Pobiera rzeczywisty adres obiektu nawet przy przeciążonym operatorze \texttt{\&}. \\
\bottomrule
\end{tabularx}
\end{center}

\subsection{Deklarowanie wskaźników}

Operator \texttt{\&} zastosowany do obiektu daje jego adres, natomiast
operator jednoargumentowy \texttt{*} służy do dereferencji (dostępu do
wskazywanego obiektu):

\begin{lstlisting}
#include <iostream>

int main() {
    int x = 10;
    int* p = &x;             // p points to x

    std::cout << x << '\n';  // 10
    std::cout << *p << '\n'; // 10

    *p = 25;                 // modifies x
    std::cout << x << '\n';  // 25
}
\end{lstlisting}

W deklaracji \texttt{int* p} gwiazdka należy do deklaratora, nie jest
operacją dereferencji. Nie należy też utożsamiać wskaźnika z samym
obiektem: \texttt{p} jest osobnym obiektem o własnej wartości i adresie.

\begin{lstlisting}
int x = 7;
int* p = &x;
int** pp = &p;  // pointer to a pointer

**pp = 9;       // changes x
\end{lstlisting}

Zapis \texttt{int* a, b;} deklaruje wskaźnik \texttt{a} oraz
zwykłą zmienną \texttt{b}. Dla czytelności lepiej deklarować je
oddzielnie.

\subsection{Wskaźnik pusty i poprawna dereferencja}

Pusty wskaźnik zapisujemy za pomocą \texttt{nullptr} (od C++11).
Nie wolno go dereferencjonować. Samo porównanie z
\texttt{nullptr} jest bezpieczne:

\begin{lstlisting}
int* p = nullptr;

if (p != nullptr) {
    *p = 5;
}

if (p) {                 // equivalent non-null test
    *p = 6;
}
\end{lstlisting}

\texttt{nullptr} ma osobny typ \texttt{std::nullptr\_t}; w odróżnieniu od
historycznych zapisów \texttt{0} i \texttt{NULL} nie jest zwykłą
liczbą całkowitą. Pomaga to przy przeciążaniu funkcji:

\begin{lstlisting}
void f(int);
void f(int*);

f(0);         // calls f(int)
f(nullptr);   // calls f(int*)
\end{lstlisting}

\paragraph{Uwaga o wskaźnikach niezainicjalizowanych.}
Lokalne \texttt{int* p;} bez inicjalizacji ma nieokreśloną wartość;
korzystanie z niej może prowadzić do niezdefiniowanego zachowania.
Zawsze inicjalizuj wskaźnik poprawnym adresem lub \texttt{nullptr}.

\subsection{Dostęp do składowych: \texttt{.} i \texttt{->}}

Operator \texttt{.} wybiera składową obiektu, natomiast
\texttt{->} wybiera składową przez wskaźnik:

\begin{lstlisting}
#include <string>

struct User {
    std::string name;
    int score = 0;
};

User u{"Anna", 12};
User* p = &u;

u.score = 13;
p->score = 14;       // same as (*p).score = 14
\end{lstlisting}

Nawiasy w \texttt{(*p).score} są potrzebne ze względu na priorytety
operatorów. Zapis \texttt{*p.score} nie oznacza tego samego.

\section{Wskaźniki, referencje i kwalifikator \texttt{const}}

\subsection{Kiedy wskaźnik, a kiedy referencja?}

Referencja \texttt{T\&} jest zwykle odpowiednia, gdy obiekt musi
istnieć i parametr nie jest opcjonalny. Wskaźnik \texttt{T*} nadaje
się, gdy brak obiektu jest dopuszczalny lub gdy chcemy zmieniać
obiekt, na który wskazujemy.

\begin{lstlisting}
void increase(int& x) { ++x; }

void increase_if_present(int* x) {
    if (x) { ++*x; }
}

int n = 3;
increase(n);             // n == 4
increase_if_present(&n); // n == 5
increase_if_present(nullptr); // no operation
\end{lstlisting}

Referencja również może stać się \emph{wisząca}, gdy obiekt,
do którego się odnosi, przestanie istnieć. Typ \texttt{T\&}
nie gwarantuje samodzielnie, że czas życia obiektu jest wystarczająco długi.

\subsection{Co jest stałe: wskaźnik czy obiekt?}

\begin{lstlisting}
int a = 1;
int b = 2;

const int* p = &a;       // read-only access through p
p = &b;                 // OK: p can point elsewhere
// *p = 3;              // error

int* const q = &a;       // q cannot be reassigned
*q = 3;                 // OK: a is mutable
// q = &b;              // error

const int* const r = &a; // neither reassignment nor write
int const* s = &a;       // same type as const int*
\end{lstlisting}

Reguła praktyczna: \texttt{const} po \texttt{*} dotyczy samego
wskaźnika, a przed \texttt{*} --- wskazywanego typu.

\paragraph{Pozorna niemutowalność.}
Wskaźnik \texttt{const T*} nie pozwala modyfikować obiektu
\emph{przez ten wskaźnik}. Nie oznacza, że obiekt nie może ulec zmianie
inną drogą:

\begin{lstlisting}
int x = 10;
const int* read_only = &x;
x = 20;                 // allowed
// *read_only = 30;     // not allowed
\end{lstlisting}

\section{Czas życia obiektu i wskaźniki wiszące}

\subsection{Adres nie oznacza prawa dostępu}

Wskaźnik może przechowywać adres miejsca, w którym obiekt już
nie istnieje. Jest wtedy wskaźnikiem \emph{wiszącym}
(ang. \emph{dangling pointer}). Dereferencja takiego wskaźnika
prowadzi do niezdefiniowanego zachowania (\emph{undefined behavior}, UB).

\begin{lstlisting}
int* bad() {
    int local = 42;
    return &local;       // WRONG: local dies at return
}

// int x = *bad();       // UB: dangling pointer
\end{lstlisting}

Prawidłowo można zwrócić samą wartość:

\begin{lstlisting}
int good() {
    int local = 42;
    return local;        // safe value return
}
\end{lstlisting}

\subsection{Unieważnianie wskaźników przez kontenery}

Nie każdy wskaźnik staje się wiszący przez wyjście z funkcji.
Relokacja elementów \texttt{std::vector} także może unieważnić
wskaźniki, referencje i iteratory do tych elementów:

\begin{lstlisting}
#include <vector>

std::vector<int> values{10, 20};
int* p = &values[0];

values.reserve(100);     // may reallocate storage
// *p = 50;              // unsafe if reallocation occurred
p = &values[0];          // reacquire a valid pointer
\end{lstlisting}

Po zmianie pojemności wymagającej realokacji adresy elementów
się zmieniają. W ogólności po operacji mogącej unieważnić wskaźnik
należy pobrać go ponownie, a nie próbować odgadnąć, czy nadal działa.

\subsection{Trzy niezależne pytania}

Przy każdym wskaźniku rozstrzygaj:
\begin{enumerate}
    \item Czy wartość wskaźnika jest poprawna (np. nie jest
    niezainicjalizowana)?
    \item Czy obiekt, na który wskazujemy, nadal istnieje?
    \item Czy operacja jest dozwolona dla danego typu i zakresu pamięci
    (np. czy nie wychodzi poza tablicę)?
\end{enumerate}

Sam test \texttt{p != nullptr} rozwiązuje co najwyżej fragment
pierwszego problemu, ale \emph{nie} wykrywa wiszących wskaźników.

\section{Tablice, arytmetyka wskaźników i \texttt{std::span}}

\subsection{Arytmetyka wskaźników}

Dla wskaźnika do elementu tablicy operacje \texttt{p + k} i
\texttt{p[k]} odnoszą się do jej elementów, pod warunkiem że
poruszamy się w dozwolonym zakresie. Można utworzyć wskaźnik
\emph{jeden za końcem} tablicy, ale nie wolno go dereferencjonować.

\begin{lstlisting}
int a[]{4, 8, 12};
int* p = a;          // points to a[0]

int x = *(p + 1);    // 8
int y = p[2];       // 12
int* end = a + 3;   // valid one-past-the-end pointer
// int z = *end;     // UB
\end{lstlisting}

Dodawanie i odejmowanie wskaźników ma sens w obrębie \emph{tej samej}
tablicy (z uwzględnieniem pozycji jeden za końcem). Nie należy
wykorzystywać takich działań do przemieszczania się między dowolnymi
obiektami w pamięci.

\subsection{Tablica to nie wskaźnik}

Typ \texttt{int[3]} i typ \texttt{int*} są różne, choć w wielu
wyrażeniach tablica jest automatycznie konwertowana na wskaźnik
do pierwszego elementu (tzw. \emph{array-to-pointer decay}).

\begin{lstlisting}
int a[3] = {1, 2, 3};
int* p = a;

static_assert(sizeof(a) == 3 * sizeof(int));
// sizeof(p) is the size of a pointer, not of the array
\end{lstlisting}

Funkcja \texttt{void f(int* p)} nie zna długości tablicy.
Jeśli trzeba pracować z ciągłym fragmentem pamięci o znanym
rozmiarze, C++20 udostępnia \texttt{std::span}:

\begin{lstlisting}
#include <span>
#include <vector>

int sum(std::span<const int> xs) {
    int result = 0;
    for (int x : xs) {
        result += x;
    }
    return result;
}

int a[]{1, 2, 3};
std::vector<int> v{4, 5, 6};

int s1 = sum(a);  // 6
int s2 = sum(v);  // 15
\end{lstlisting}

\texttt{std::span<T>} przechowuje widok ciągłych danych
(z informacją o liczbie elementów), ale \emph{nie jest właścicielem}
pamięci. Gdy dane znikną albo przeniosą się w inne miejsce,
\texttt{span} również może stać się niepoprawny. Dla większości
dynamicznych tablic własnościowych preferuj \texttt{std::vector<T>}.

\section{\texttt{void*}, konwersje i rzutowania wskaźników}

\subsection{\texttt{void*}: adres bez informacji o typie}

Wskaźnik \texttt{void*} może przechowywać adres obiektu bez
informacji o jego konkretnym typie. Bywa używany w interfejsach C
i bibliotekach niskopoziomowych, ale nie pozwala na bezpośrednią
dereferencję:

\begin{lstlisting}
int x = 42;
void* opaque = &x;              // type erased
// *opaque = 7;                 // error: cannot dereference void*

int* p = static_cast<int*>(opaque);
*p = 7;                        // valid: original object is int
\end{lstlisting}

Przywrócenie błędnego typu i późniejsza próba użycia obiektu
mogą prowadzić do niezdefiniowanego zachowania. Jeśli typ jest
znany, zwykle lepiej przechowywać od razu \texttt{T*}.

\subsection{\texttt{static\_cast}, \texttt{dynamic\_cast} i
\texttt{reinterpret\_cast}}

\begin{itemize}
    \item \texttt{static\_cast<T*>} nadaje się m.in. do poprawnych
    konwersji pomiędzy \texttt{void*} a wskaźnikiem obiektowym,
    jeśli rzeczywisty typ i wyrównanie na to pozwalają.
    \item \texttt{dynamic\_cast<T*>} służy do sprawdzanych w czasie
    wykonania konwersji w hierarchiach klas polimorficznych.
    Niepowodzenie konwersji wskaźnika zwraca \texttt{nullptr}.
    \item \texttt{reinterpret\_cast<T*>} umożliwia niskopoziomową
    reinterpretację reprezentacji wskaźnika, ale sama konwersja
    nie uprawnia do dowolnej dereferencji ani nie tworzy obiektu
    odpowiedniego typu.
\end{itemize}

\begin{lstlisting}
struct Base {
    virtual ~Base() = default;
};

struct Derived : Base {
    int value = 12;
};

Derived d;
Base* base = &d;

if (Derived* p = dynamic_cast<Derived*>(base)) {
    p->value = 30;          // checked downcast succeeded
}
\end{lstlisting}

Przy przekazywaniu parametrów funkcji nie używaj rzutowań
jako sposobu na obejście \texttt{const}, sprawdzania czasu życia
ani zasad własności. Dobrze zaprojektowany interfejs zazwyczaj
w ogóle nie wymaga rzutowania wskaźników.

\section{Pamięć dynamiczna: \texttt{new}, \texttt{delete} i ich ograniczenia}

\subsection{Klasyczny mechanizm}

Operator \texttt{new} tworzy obiekt w pamięci przydzielonej
dynamicznie i zwraca wskaźnik. Odpowiadający mu
\texttt{delete} niszczy obiekt i zwalnia pamięć:

\begin{lstlisting}
int* p = new int(42);
// use *p

delete p;
p = nullptr;           // avoids this one stale alias
\end{lstlisting}

Dla tablic obowiązuje para \texttt{new[]} i \texttt{delete[]}:

\begin{lstlisting}
int* a = new int[5]{}; // zero-initialized elements
// use a[0], ..., a[4]

delete[] a;
\end{lstlisting}

\textbf{Nie wolno} łączyć \texttt{new[]} z \texttt{delete} ani
zwykłego \texttt{new} z \texttt{delete[]}. Ponowne
\texttt{delete} tego samego żywego przydziału również jest błędem.
\texttt{delete nullptr;} jest natomiast bezpieczne.

\subsection{Dlaczego ręczne zarządzanie pamięcią jest niebezpieczne?}

\begin{lstlisting}
void fragile() {
    int* p = new int(42);
    operation_that_may_throw();
    delete p;   // skipped if the operation throws
}
\end{lstlisting}

Jeśli funkcja zgłosi wyjątek, pamięć zostaje utracona
(\emph{memory leak}). Analogiczny problem powstaje przy
wielu instrukcjach \texttt{return}, kopiowaniu wskaźników i
niespójnym przekazywaniu odpowiedzialności za \texttt{delete}.

\paragraph{Zalecenie.}
W typowym kodzie aplikacyjnym nie pisz bezpośrednio
\texttt{new}/\texttt{delete}. Użyj obiektu lokalnego,
\texttt{std::vector<T>} albo inteligentnego wskaźnika.
Ręczne operacje alokowania pozostają potrzebne w implementacjach
niskopoziomowych, alokatorach i przy integracji z określonym API.

\section{RAII i semantyka własności}

\subsection{RAII: zasób należy do obiektu}

\emph{Resource Acquisition Is Initialization} (RAII) oznacza
powiązanie przejęcia zasobu z utworzeniem obiektu zarządzającego,
a zwolnienia zasobu z jego destruktorem. Stosuje się je nie tylko
do pamięci, ale też do plików, blokad, gniazd sieciowych i uchwytów
systemowych.

\begin{lstlisting}
#include <fstream>
#include <string>

void save(const std::string& path) {
    std::ofstream file(path);
    file << "example\n";
} // file is closed by its destructor
\end{lstlisting}

Destruktor uruchamia się podczas normalnego opuszczania zakresu
oraz podczas odwijania stosu po wyjątku. Stąd RAII upraszcza
bezpieczeństwo wyjątkowe.

\subsection{Trzy rodzaje relacji z obiektem}

\begin{description}
    \item[Własność wyłączna (\emph{unique ownership}).]
    Dokładnie jeden obiekt zarządzający odpowiada za zniszczenie
    zasobu: najczęściej \texttt{std::unique\_ptr<T>}.
    \item[Własność współdzielona (\emph{shared ownership}).]
    Wielu właścicieli może utrzymywać jeden zasób przy życiu:
    \texttt{std::shared\_ptr<T>}.
    \item[Dostęp bez własności (\emph{non-owning access}).]
    Użytkownik widzi obiekt, ale nie przedłuża jego życia:
    zwykły wskaźnik, referencja, \texttt{std::span<T>} lub
    \texttt{std::weak\_ptr<T>} (ten ostatni obserwuje stan
    zarządzany przez \texttt{shared\_ptr} i pozwala warunkowo
    uzyskać właściciela).
\end{description}

Dobór reprezentacji wynika z \emph{kontraktu własności}, a nie
z samego faktu, że obiekt znajduje się w pamięci dynamicznej.

\section{\texttt{std::unique\_ptr}: jeden właściciel}

\begin{center}\small
\begin{tabularx}{\textwidth}{@{}p{0.36\textwidth}X@{}}
\toprule
Mechanizm & Własność i czas życia \\
\midrule
\code{std::make_unique<T>(...)} & Tworzy obiekt i zwraca jego jedynego właściciela. \\
\code{unique_ptr::get()} & Zwraca surowy wskaźnik bez przekazania własności. \\
\code{unique_ptr::reset(p)} & Niszczy obecnie posiadany obiekt i przejmuje wskaźnik \code{p}. \\
\code{unique_ptr::release()} & Oddaje surowy wskaźnik bez zniszczenia obiektu; odtąd trzeba osobno zadbać o jego zwolnienie. \\
\code{std::shared_ptr<T>} & Współdzieli własność; obiekt znika po odejściu ostatniego właściciela. \\
\code{weak_ptr::lock()} & Próbuje uzyskać \code{shared_ptr}; zwraca pusty wynik, jeśli obiekt wygasł. \\
\bottomrule
\end{tabularx}
\end{center}

\subsection{Tworzenie obiektu i automatyczne niszczenie}

\texttt{std::unique\_ptr<T>} (nagłówek \texttt{<memory>}, C++11)
zarządza pojedynczym obiektem przez wyłączną własność. Od C++14
najczęściej tworzymy go przez \texttt{std::make\_unique<T>}:

\begin{lstlisting}
#include <memory>
#include <string>

#include <utility>

struct Person {
    explicit Person(std::string n) : name(std::move(n)) {}
    std::string name;
};

void example() {
    auto p = std::make_unique<Person>("Ala");
    p->name = "Ola";
    // no delete needed
} // destructor destroys Person
\end{lstlisting}

\texttt{unique\_ptr} można traktować jako obiekt zarządzający:
jego destruktor wywołuje odpowiednią funkcję usuwającą zasób
(\emph{deleter}). Domyślnie jest nią \texttt{std::default\_delete<T>}.

\subsection{Brak kopiowania i przenoszenie własności}

Wskaźnik o własności wyłącznej nie może być kopiowany.
Może natomiast zostać \emph{przeniesiony} za pomocą
\texttt{std::move}:

\begin{lstlisting}
#include <memory>
#include <utility>

auto first = std::make_unique<int>(10);
// auto bad = first;                  // error: copy

auto second = std::move(first);      // transfers ownership
// first == nullptr
*second = 20;
\end{lstlisting}

\texttt{std::move} samo niczego nie zwalnia i nie przenosi fizycznie
wskazywanego obiektu. Zwraca wyrażenie pozwalające wywołać konstruktor
lub przypisanie przenoszące. W tym przypadku przenoszona jest
\emph{własność wskaźnika}, a obiekt typu \texttt{int} nie zmienia adresu.

\subsection{Przekazywanie do funkcji: kto jest właścicielem?}

Rozróżnijmy trzy interfejsy:

\begin{lstlisting}
#include <memory>
#include <utility>

struct Widget { void run() {} };

void inspect(const Widget& w);             // borrows, non-null
void maybe_inspect(const Widget* w);       // borrows, nullable
void take(std::unique_ptr<Widget> w);      // takes ownership

void demo() {
    auto p = std::make_unique<Widget>();
    inspect(*p);                // p still owns the object
    maybe_inspect(p.get());     // p still owns the object
    take(std::move(p));         // ownership transferred
    // p is now empty
}
\end{lstlisting}

Jeżeli funkcja \emph{jedynie używa} obiektu, przekazywanie przez
\texttt{T\&}, \texttt{const T\&} lub \texttt{T*} jest zwykle
czytelniejsze niż przekazywanie \texttt{unique\_ptr} przez referencję.
Parametr \texttt{unique\_ptr<T>} przyjmowany \emph{przez wartość}
sygnalizuje przejęcie własności. Dla funkcji, która ma zastąpić
obiekt będący własnością wywołującego, można użyć
\texttt{std::unique\_ptr<T>\&}.

\subsection{Metody \texttt{get}, \texttt{reset}, \texttt{release}, \texttt{swap}}

\begin{lstlisting}
#include <memory>
#include <utility>

auto p = std::make_unique<int>(7);
int* observer = p.get();       // borrow; do not delete

p.reset(new int(9));           // destroys old int
// observer is now dangling

auto q = std::make_unique<int>(11);
p.swap(q);                    // exchanges managed objects

int* raw = q.release();        // q becomes empty; raw owns now
// responsibility for delete raw is now explicit
delete raw;
\end{lstlisting}

\texttt{get()} zwraca dostęp bez własności. \texttt{reset()} zastępuje
zarządzany obiekt, zwalniając dotychczasowy. \texttt{release()} oddaje
surowy wskaźnik \emph{bez zwolnienia obiektu}; dlatego jest potrzebne
głównie przy przekazywaniu własności do interfejsu, który tego wymaga.
Nie należy mylić \texttt{release()} z \texttt{reset()}.

\subsection{Tablice: \texttt{unique\_ptr<T[]>} a \texttt{vector<T>}}

\begin{lstlisting}
#include <memory>
#include <vector>

auto a = std::make_unique<int[]>(5); // array of 5 zeros
a[2] = 17;

std::vector<int> b(5);         // usually more convenient
b.at(2) = 17;                  // checked indexing
\end{lstlisting}

\texttt{unique\_ptr<T[]>} dobiera \texttt{delete[]} i udostępnia
\texttt{operator[]}, lecz sam nie przechowuje rozmiaru.
Zwykle preferuj \texttt{std::vector<T>}, bo posiada rozmiar,
iteratory i bogatszy interfejs. Dla tablicy o stałym rozmiarze
określonym w kodzie rozważ \texttt{std::array<T, N>}.

\subsection{Własny deleter: zasób inny niż pamięć \texttt{new}}

Inteligentny wskaźnik może wywoływać dowolną funkcję zwalniającą
zasób. Typowym przykładem jest plik w interfejsie C:

\begin{lstlisting}
#include <cstdio>
#include <memory>

struct FileCloser {
    void operator()(std::FILE* file) const noexcept {
        if (file) { std::fclose(file); }
    }
};

using FilePtr = std::unique_ptr<std::FILE, FileCloser>;

void example() {
    FilePtr file(std::fopen("log.txt", "w"));
    if (!file) { return; }

    std::fputs("hello\n", file.get());
} // fclose is called automatically
\end{lstlisting}

W praktyce deleter bywa osobną strukturą, a jego typ wchodzi
w skład typu \texttt{unique\_ptr<T, Deleter>}. Mechanizm pozwala
owinąć RAII również zasoby zwalniane przez \texttt{free()},
\texttt{fclose()} lub bibliotekę zewnętrzną.

\subsection{Polimorfizm i wirtualny destruktor}

\begin{lstlisting}
#include <memory>

struct Shape {
    virtual ~Shape() = default;
    virtual double area() const = 0;
};

struct Circle : Shape {
    explicit Circle(double r) : radius(r) {}
    double area() const override { return 3.14159 * radius * radius; }
    double radius;
};

std::unique_ptr<Shape> make_shape() {
    return std::make_unique<Circle>(2.0);
}
\end{lstlisting}

Przy niszczeniu obiektu pochodnego przez wskaźnik do bazy
zwykle potrzebny jest wirtualny destruktor bazy.
Dotyczy to również przenoszenia własności z klasy pochodnej
do bazowej przy użyciu \texttt{unique\_ptr} i domyślnego deletera.

\section{\texttt{std::shared\_ptr}: współdzielona własność}

\subsection{Jak działa licznik właścicieli?}

\texttt{std::shared\_ptr<T>} (C++11) udostępnia współdzieloną
własność zasobu. Zwykle tworzymy go za pomocą
\texttt{std::make\_shared<T>}:

\begin{lstlisting}
#include <memory>

auto a = std::make_shared<int>(10);
auto b = a;             // shared ownership

// a.use_count() == 2 (provided there are no other owners)
b.reset();              // releases b's ownership
// a.use_count() == 1
\end{lstlisting}

Gdy liczba \emph{silnych właścicieli} spadnie do zera, zarządzany
obiekt zostaje zniszczony. Implementacja utrzymuje tzw.
\emph{blok kontrolny} (\emph{control block}) zawierający między
innymi liczniki właścicieli, informacje potrzebne do zniszczenia
zasobu i ewentualny deleter.

\paragraph{Schemat ideowy.}
\begin{lstlisting}
shared_ptr a ---\
                +--> [ control block: strong = 2 ] --> [ T ]
shared_ptr b ---/           weak = 0
\end{lstlisting}

\texttt{use\_count()} ma znaczenie przede wszystkim diagnostyczne.
Nie powinno się na podstawie odczytanego licznika podejmować decyzji
o bezpieczeństwie wątkowym ani o wyłącznej własności, która może
zmienić się równolegle.

\subsection{Tworzenie i przekazywanie \texttt{shared\_ptr}}

\begin{lstlisting}
#include <memory>

struct Data { int x = 0; };

std::shared_ptr<Data> create() {
    return std::make_shared<Data>();
}

void keep(std::shared_ptr<Data> p) {
    // stores a copy of p somewhere, if needed
}

void read_only(const Data& data) {
    // does not require shared ownership
}

void demo() {
    auto p = create();
    read_only(*p); // borrow
    keep(p);       // creates another temporary shared owner
}
\end{lstlisting}

Parametr \texttt{shared\_ptr<T>} przyjmowany przez wartość oznacza,
że funkcja otrzymuje własny udział we własności na czas wywołania
i ewentualnie może go przechować. Jeśli funkcja wyłącznie
czyta obiekt, zwykle wystarcza \texttt{const T\&} lub \texttt{T*}.
Samo przekazanie \texttt{const shared\_ptr<T>\&} unika inkrementacji
licznika, ale bardziej wiąże interfejs z modelem własności.

\subsection{\texttt{make\_shared} a konstruktor ze wskaźnika}

W typowych implementacjach \texttt{make\_shared} umieszcza obiekt
i blok kontrolny w jednym przydziale pamięci, co może ograniczać
narzut alokacji. Natomiast konstrukcja \texttt{shared\_ptr<T>(new T)}
zwykle wymaga osobnych przydziałów.

\begin{lstlisting}
auto preferred = std::make_shared<int>(42);
std::shared_ptr<int> also_valid(new int(42));
\end{lstlisting}

Nie wolno jednak tworzyć dwóch \texttt{shared\_ptr} niezależnie
z \emph{tego samego} surowego wskaźnika:

\begin{lstlisting}
int* raw = new int(42);
std::shared_ptr<int> a(raw);
// std::shared_ptr<int> b(raw); // WRONG: second control block
\end{lstlisting}

Dwa niezależne bloki kontrolne próbowałyby zwolnić ten sam obiekt.
Poprawne współdzielenie powstaje przez \emph{kopiowanie istniejącego}
\texttt{shared\_ptr}: \texttt{auto b = a;}.

\paragraph{Kiedy nie stosować \texttt{make\_shared}?}
Zintegrowany przydział może zatrzymywać pamięć zajmowaną przez
obiekt w przydzielonym bloku do momentu zniszczenia także ostatniego
\texttt{weak\_ptr}, nawet gdy destruktor obiektu już się wykonał.
Przy dużych obiektach i długo żyjących słabych referencjach
warto rozważyć osobny przydział. Innymi powodami mogą być
specjalne wymagania alokacji lub konstruktory niedostępne
dla \texttt{make\_shared}.

\subsection{Współdzielenie własności a współdzielenie mutacji}

Sam \texttt{shared\_ptr} nie synchronizuje dostępu do pól obiektu:

\begin{lstlisting}
#include <memory>
#include <mutex>

struct Counter {
    std::mutex mutex;
    int value = 0;

    void increment() {
        std::lock_guard<std::mutex> lock(mutex);
        ++value;
    }
};

auto p = std::make_shared<Counter>();
// many threads may own copies of p, but must synchronize
// concurrent modifications of Counter::value
\end{lstlisting}

Blok kontrolny pozwala różnym egzemplarzom \texttt{shared\_ptr}
bezpiecznie współdzielić własność. Równoczesna modyfikacja
\emph{tego samego} egzemplarza wymaga jednak synchronizacji.
Od C++20 można użyć \texttt{std::atomic<std::shared\_ptr<T>>}.

\section{\texttt{std::weak\_ptr}: obserwacja bez utrzymywania przy życiu}

\subsection{Problem cykli własności}

Gdy obiekty wzajemnie posiadają się przez \texttt{shared\_ptr},
liczniki właścicieli mogą nigdy nie spaść do zera:

\begin{lstlisting}
#include <memory>

struct Node {
    std::shared_ptr<Node> next;
};

void leak_cycle() {
    auto a = std::make_shared<Node>();
    auto b = std::make_shared<Node>();
    a->next = b;
    b->next = a;          // cycle of strong ownership
} // nodes remain alive after both locals are destroyed
\end{lstlisting}

To nie jest problem \emph{niezerowego wskaźnika}, lecz błędnego
modelu własności: każdy węzeł utrzymuje przy życiu drugi.

\subsection{Przerywanie cyklu i \texttt{lock()}}

\texttt{std::weak\_ptr<T>} obserwuje ten sam blok kontrolny,
ale nie zwiększa liczby silnych właścicieli:

\begin{lstlisting}
#include <memory>

struct Node {
    std::weak_ptr<Node> parent;         // non-owning
    std::shared_ptr<Node> child;        // owning
};

void example() {
    auto root = std::make_shared<Node>();
    auto leaf = std::make_shared<Node>();

    root->child = leaf;
    leaf->parent = root;

    if (auto parent = leaf->parent.lock()) {
        // parent is a temporary shared owner
        (void)parent;
    }
} // no ownership cycle
\end{lstlisting}

Metoda \texttt{lock()} atomowo próbuje uzyskać
\texttt{shared\_ptr<T>}. Jeżeli obiekt jeszcze istnieje,
zwraca właściciela utrzymującego go przy życiu przez czas
używania. W przeciwnym razie zwraca pusty \texttt{shared\_ptr}.

\begin{lstlisting}
std::weak_ptr<int> observer;
{
    auto p = std::make_shared<int>(8);
    observer = p;
    if (auto alive = observer.lock()) {
        *alive += 1;           // safe while alive owns it
    }
}

// observer.expired() is now true
// observer.lock() now returns an empty shared_ptr
\end{lstlisting}

W kodzie współbieżnym nie stosuj schematu
\texttt{if (!w.expired()) { ... }} z późniejszym odwołaniem:
między sprawdzeniem a użyciem obiekt może zniknąć.
Preferuj jedną operację \texttt{lock()} i posługuj się
otrzymanym \texttt{shared\_ptr}.

\section{\texttt{enable\_shared\_from\_this} i wskaźnik \texttt{this}}

Gdy metoda obiektu zarządzanego przez \texttt{shared\_ptr}
chce przekazać dalej \emph{współdzieloną własność do siebie},
nie może wykonać \texttt{std::shared\_ptr<T>(this)}:
powstałby drugi, niezależny blok kontrolny.

Można zamiast tego dziedziczyć publicznie po
\texttt{std::enable\_shared\_from\_this<T>}:

\begin{lstlisting}
#include <memory>

struct Session : std::enable_shared_from_this<Session> {
    std::shared_ptr<Session> self() {
        return shared_from_this();
    }
};

void example() {
    auto session = std::make_shared<Session>();
    auto second_owner = session->self();
    // both share the same control block
}
\end{lstlisting}

Metodę \texttt{shared\_from\_this()} wolno wywołać dopiero wtedy,
gdy obiektem zarządza już \texttt{shared\_ptr}. Jej wywołanie
na obiekcie lokalnym albo w konstruktorze, zanim powstanie
właściciel, zgłasza wyjątek \texttt{std::bad\_weak\_ptr}.

Od C++17 dostępna jest również metoda \texttt{weak\_from\_this()}.

\section{Zaawansowany przykład: wskaźnik aliasujący \texttt{shared\_ptr}}

\texttt{shared\_ptr} rozróżnia \emph{obiekt zarządzany} (którego
czas życia kontroluje blok kontrolny) i \emph{przechowywany
wskaźnik} (zwracany przez \texttt{get()}). Specjalny konstruktor
aliasujący pozwala udostępnić adres składowej, jednocześnie
utrzymując przy życiu cały obiekt:

\begin{lstlisting}
#include <memory>

struct Pair {
    int first = 1;
    int second = 2;
};

void example() {
    auto whole = std::make_shared<Pair>();
    std::shared_ptr<int> member(whole, &whole->second);

    whole.reset();       // Pair is still alive
    *member = 20;        // accesses Pair::second safely
} // member destroys the last strong owner
\end{lstlisting}

Konstruktor aliasujący jest przydatny np. przy udostępnianiu
podobiektów, ale wymaga prawidłowej relacji czasów życia.
Nie należy go traktować jako sposobu na naprawienie
wcześniej wiszącego wskaźnika.

\section{Wskaźniki w interfejsach funkcji i klas}

\subsection{Katalog typowych kontraktów}

\begin{center}
\small
\begin{tabular}{|p{0.31\textwidth}|p{0.56\textwidth}|}
\hline
\textbf{Interfejs} & \textbf{Znaczenie} \\
\hline
\texttt{void f(const T\&)} & Obiekt wymagany, tylko odczyt; brak przekazania własności. \\
\hline
\texttt{void f(T\&)} & Obiekt wymagany, możliwa modyfikacja; brak przekazania własności. \\
\hline
\texttt{void f(const T*)} & Obiekt opcjonalny, tylko odczyt; brak własności. \\
\hline
\texttt{void f(T*)} & Obiekt opcjonalny, możliwa modyfikacja; brak własności. \\
\hline
\texttt{void f(unique\_ptr<T>)} & Przejęcie wyłącznej własności. \\
\hline
\texttt{unique\_ptr<T> f()} & Przekazanie wyłącznej własności wywołującemu. \\
\hline
\texttt{void f(shared\_ptr<T>)} & Uzyskanie udziału we współdzielonej własności. \\
\hline
\texttt{void f(weak\_ptr<T>)} & Możliwość obserwacji obiektu, którego życie kontrolują inni. \\
\hline
\texttt{void f(span<const T>)} & Widok na ciągłą sekwencję, bez własności. \\
\hline
\end{tabular}
\end{center}

W tabeli pominięto przestrzeń nazw \texttt{std::} dla czytelności.
To \emph{konwencje projektowania API}, a nie automatyczne
gwarancje długowieczności surowych wskaźników i referencji.

\subsection{Dlaczego nie zawsze używać \texttt{shared\_ptr}?}

Rozważmy funkcję, która wyświetla nazwę:

\begin{lstlisting}
#include <iostream>
#include <memory>
#include <string>

struct Product { std::string name; };

void print(const Product& p) {
    std::cout << p.name << '\n';
}

void example() {
    Product local{"book"};
    print(local);                  // no allocation

    auto owned = std::make_unique<Product>(Product{"pen"});
    print(*owned);                 // still no extra owner
}
\end{lstlisting}

Funkcja \texttt{print} nie potrzebuje znać strategii
przechowywania obiektu. Przekazanie referencji oddziela
\emph{sposób dostępu} od \emph{sposobu zarządzania czasem życia}.

\section{PImpl: ukrywanie implementacji przez \texttt{unique\_ptr}}

Technika \emph{pointer to implementation} pozwala utrzymać
szczegóły implementacji poza publicznym nagłówkiem.
Typ implementacyjny może być niekompletny w nagłówku, ale
miejsce wykonania domyślnego deletera musi widzieć jego definicję.
Dlatego destruktor klasy definiujemy w pliku \texttt{.cpp}.

\begin{lstlisting}
// widget.hpp
#pragma once
#include <memory>

class Widget {
public:
    Widget();
    ~Widget();                 // defined in .cpp

    Widget(Widget&&) noexcept;
    Widget& operator=(Widget&&) noexcept;

    Widget(const Widget&) = delete;
    Widget& operator=(const Widget&) = delete;

    void work();

private:
    struct Impl;
    std::unique_ptr<Impl> impl_;
};
\end{lstlisting}

\begin{lstlisting}
// widget.cpp
#include "widget.hpp"
#include <memory>

struct Widget::Impl {
    int value = 0;
    void run() { ++value; }
};

Widget::Widget() : impl_(std::make_unique<Impl>()) {}
Widget::~Widget() = default;
Widget::Widget(Widget&&) noexcept = default;
Widget& Widget::operator=(Widget&&) noexcept = default;

void Widget::work() { impl_->run(); }
\end{lstlisting}

Ten przykład prezentuje semantykę przenoszenia. Po przeniesieniu
źródłowy obiekt \texttt{Widget} ma pusty \texttt{impl\_};
nie wolno wywoływać na nim \texttt{work()} bez odpowiedniego
sprawdzenia. W produkcyjnym API można ten stan jawnie obsłużyć
lub określić w kontrakcie.

\section{Wskaźniki do funkcji i składowych klas}

Wskaźniki nie muszą odnosić się do danych. Można wskazywać
funkcje, a także składowe klas, choć te ostatnie wymagają
innej składni.

\subsection{Wskaźnik do funkcji}

\begin{lstlisting}
int add(int a, int b) { return a + b; }
int mul(int a, int b) { return a * b; }

using Operation = int(*)(int, int);

int apply(int x, int y, Operation op) {
    return op(x, y);
}

int a = apply(2, 3, &add); // 5
int b = apply(2, 3, &mul); // 6
\end{lstlisting}

Obiekty wywoływalne różnych typów można przechowywać za pomocą
\texttt{std::function}. Prosty szablon lub lambda bywają jednak
bardziej efektywne.

\subsection{Wskaźnik do składowej}

\begin{lstlisting}
struct Record {
    int value = 0;
    int twice() const { return 2 * value; }
};

int Record::* field = &Record::value;
int (Record::* method)() const = &Record::twice;

Record r{7};
int x = r.*field;          // 7
int y = (r.*method)();     // 14

Record* p = &r;
int z = p->*field;         // 7
\end{lstlisting}

Wskaźnik do składowej nie jest zwykłym adresem gotowego obiektu.
Dopiero w połączeniu z konkretną instancją klasy pozwala
odwołać się do odpowiedniej składowej.

\section{Wskaźniki i wielowątkowość}

\subsection{\texttt{shared\_ptr} nie eliminuje wyścigów danych}

Możliwość bezpiecznego utrzymania czasu życia obiektu nie oznacza,
że jego pola można dowolnie zmieniać z wielu wątków.
Jeśli dwa wątki zapisują do \texttt{*p}, potrzebna jest osobna
synchronizacja (np. \texttt{std::mutex} albo atomiki).

Od C++20 dostępna jest specjalizacja
\texttt{std::atomic<std::shared\_ptr<T>>}, przydatna do atomowego
publikowania nowej wersji współdzielonego obiektu:

\begin{lstlisting}
#include <atomic>
#include <memory>

struct Config {
    int workers = 4;
};

std::atomic<std::shared_ptr<const Config>> current{
    std::make_shared<const Config>(Config{4})
};

void publish(int count) {
    current.store(std::make_shared<const Config>(Config{count}));
}

int read_workers() {
    auto snapshot = current.load(); // keeps snapshot alive
    return snapshot->workers;
}
\end{lstlisting}

Tu czytelnicy otrzymują stabilny, niemutowalny \emph{snapshot}.
Atomowość dotyczy odczytu i zapisu \texttt{shared\_ptr},
a nie arbitralnych operacji na jego zawartości.

\section{Najczęstsze błędy zebrane w jednym miejscu}

\begin{enumerate}
    \item \textbf{Dereferencja \texttt{nullptr} lub wskaźnika
    niezainicjalizowanego.} Zawsze sprawdzaj kontrakt i inicjalizuj
    wskaźniki.
    \item \textbf{Wskaźnik wiszący.} \texttt{get()}, \texttt{span}
    i surowa referencja nie przedłużają życia zasobu.
    \item \textbf{Podwójne \texttt{delete}.} Jeden przydział musi mieć
    jedną, jednoznaczną strategię zwalniania.
    \item \textbf{Mieszanie \texttt{new[]} i \texttt{delete}.}
    Do tablic służy \texttt{delete[]}; zwykle lepiej wybrać kontener.
    \item \textbf{Zwracanie adresu lokalnej zmiennej.}
    Jej czas życia kończy się przy wyjściu z zakresu.
    \item \textbf{Przekazywanie surowego wskaźnika z \texttt{get()}
    do nowego \texttt{shared\_ptr}.} Tworzy niezależny blok kontrolny.
    \item \textbf{Nieuzasadnione \texttt{shared\_ptr} wszędzie.}
    Współdzielenie własności kosztuje i zaciera odpowiedzialność.
    \item \textbf{Cykl silnych właścicieli.} Przerywaj go przez
    \texttt{weak\_ptr} lub przeprojektowanie relacji.
    \item \textbf{Nadużywanie \texttt{release()}.} Po wywołaniu
    ktoś musi jawnie przejąć odpowiedzialność za zasób.
    \item \textbf{Brak wirtualnego destruktora w bazie.}
    Usuwanie obiektu pochodnego przez wskaźnik bazowy z domyślnym
    deleterem może wówczas prowadzić do UB.
    \item \textbf{Założenie, że \texttt{shared\_ptr} chroni obiekt
    przed wyścigami.} Chroni zarządzanie własnością, nie samą mutację.
    \item \textbf{Przechowywanie wskaźników do elementów \texttt{vector}
    po realokacji.} Przy operacjach zmieniających bufor pobieraj
    adresy ponownie.
\end{enumerate}

\section{Jak wybrać właściwy typ?}

\begin{center}
\small
\begin{tabular}{|p{0.43\textwidth}|p{0.43\textwidth}|}
\hline
\textbf{Sytuacja} & \textbf{Najczęstszy wybór} \\
\hline
Pojedynczy obiekt o prostym czasie życia & Zwykła zmienna \texttt{T} \\
\hline
Dynamiczna lista elementów & \texttt{std::vector<T>} \\
\hline
Obiekt wymagany tylko do odczytu & \texttt{const T\&} \\
\hline
Obiekt wymagany i modyfikowany & \texttt{T\&} \\
\hline
Obiekt opcjonalny bez przekazania własności & \texttt{T*} lub \texttt{const T*} \\
\hline
Pojedynczy właściciel dynamicznego obiektu & \texttt{std::unique\_ptr<T>} \\
\hline
Kilku niezależnych właścicieli & \texttt{std::shared\_ptr<T>} \\
\hline
Obserwacja obiektu współdzielonego & \texttt{std::weak\_ptr<T>} \\
\hline
Fragment istniejącej tablicy, bez własności & \texttt{std::span<T>} \\
\hline
Polimorficzna kolekcja obiektów & \texttt{std::vector} zawierający \texttt{unique\_ptr} \\
\hline
\end{tabular}
\end{center}

\subsection{Praktyczny przykład: kolekcja obiektów polimorficznych}

\begin{lstlisting}
#include <iostream>
#include <memory>
#include <vector>

struct Task {
    virtual ~Task() = default;
    virtual void run() const = 0;
};

struct PrintTask : Task {
    void run() const override { std::cout << "hello\n"; }
};

int main() {
    std::vector<std::unique_ptr<Task>> tasks;
    tasks.push_back(std::make_unique<PrintTask>());

    for (const auto& task : tasks) {
        task->run();
    }
} // vector and tasks destroy everything automatically
\end{lstlisting}

\section{Podsumowanie i praktyczne reguły}

\begin{itemize}
    \item Preferuj \emph{przechowywanie przez wartość}, jeśli
    obiekt nie wymaga odrębnej, dynamicznej tożsamości.
    \item Gdy obiekt ma jednego właściciela, używaj
    \texttt{std::unique\_ptr}; własność przenoś przez
    \texttt{std::move}.
    \item Sięgaj po \texttt{std::shared\_ptr} dopiero wtedy,
    gdy kilka części programu \emph{rzeczywiście} musi niezależnie
    utrzymywać ten sam obiekt przy życiu.
    \item Jeśli obserwator nie powinien utrzymywać obiektu przy życiu,
    użyj \texttt{std::weak\_ptr} i metody \texttt{lock()}.
    \item Do niewłasnościowego dostępu używaj
    \texttt{T\&}, \texttt{const T\&}, \texttt{T*} lub
    \texttt{std::span<T>} w zależności od kontraktu.
    \item Nie dereferencjonuj wskaźnika tylko dlatego, że jest
    różny od \texttt{nullptr}: obiekt może już nie istnieć.
    \item Preferuj \texttt{std::vector} i \texttt{std::array}
    zamiast ręcznie alokowanych tablic.
    \item Zasoby systemowe obejmuj RAII i odpowiednim deleterem;
    osobno analizuj bezpieczeństwo wątkowe wskazywanego obiektu.
\end{itemize}

\paragraph{Konkluzja.}
Najważniejsze pytanie brzmi nie ``jak zdobyć adres?'', ale
``\emph{kto jest właścicielem obiektu, jak długo on istnieje
i kto ma prawo z niego korzystać?}''. Właśnie te trzy kwestie
rozstrzygają, czy w danym miejscu potrzebujemy referencji,
surowego wskaźnika, \texttt{unique\_ptr}, \texttt{shared\_ptr}
czy \texttt{weak\_ptr}.

% Dalsza lektura:
% https://en.cppreference.com/cpp/memory/unique_ptr
% https://en.cppreference.com/cpp/memory/shared_ptr
% https://en.cppreference.com/cpp/memory/weak_ptr
% https://en.cppreference.com/cpp/memory/enable_shared_from_this
% https://en.cppreference.com/cpp/container/span
